% Preserved full-split certificate development ledger.
% Source: localCodeStructured.tex before the compact main-spine rewrite.
% This file is not included in main_permConv.tex.

% Auto-extracted during the research cleanup pass. Original source is preserved unchanged.
% Main-spine structured local-code construction and finite full-split certificate.

\section{Structured local-code construction}
\label{sec:structured-local-code}

\subsection{Local-code outer certificate interface}
\label{subsec:local-code-outer-interface}

The structured construction replaces the sliding random-banded outer by a direct sum of short local codes.  The proof
interface is intentionally local-code independent: the outer contributes only its weight enumerator, while the inner
contributes slice bounds for interleaved words of a given weight.  Systematic form is not required for the outer code.

\begin{lemma}[Block-outer first-moment interface]
\label{lem:block-outer-first-moment-interface}
Let \(C_{\mathrm{loc}}\subseteq\F_2^{n_0}\) be a binary linear \([n_0,b]\) code with homogeneous weight enumerator
\[
W_{\mathrm{loc}}(z)=\sum_j A^{\mathrm{loc}}_j z^j,
\]
and let \(C_{\mathrm{out}}=C_{\mathrm{loc}}^{\oplus M}\).  Its message length is \(k=Mb\), its
outer/interleaver length is \(N=Mn_0\), and its weight enumerator is
\[
W_{\mathrm{out}}(z)=W_{\mathrm{loc}}(z)^M,
\]
where \(A_0^{\mathrm{out}}=1\) is omitted in the nonzero-codeword sum.  If \(p_h^{\mathrm{in}}(\delta)\) is any valid
upper bound on the fixed-tap dense recursive inner producing output weight at most \(\lfloor\delta N\rfloor\) from an
interleaved outer word of weight \(h\), then
\[
\Pr[d_{\min}(\Csc)\le \lfloor\delta N\rfloor]
\le
\mu_{\mathrm{block}}(\delta)
:=
\sum_{h=1}^{N} A_h^{\mathrm{out}}p_h^{\mathrm{in}}(\delta).
\]
\end{lemma}

\begin{proof}
The direct-sum enumerator identity follows by multiplying the local generating functions block by block.  Conditional
on any nonzero outer word of weight \(h\), the uniform interleaver randomizes only the support positions seen by the
inner code, so the same inner bound \(p_h^{\mathrm{in}}(\delta)\) applies to every such word.  Summing these failure
probabilities over all nonzero outer words gives \(\mu_{\mathrm{block}}(\delta)\), and Markov's inequality gives the
minimum-distance bound.
\end{proof}

\begin{lemma}[Finite block-certificate template]
\label{lem:finite-block-certificate-template}
In the setting of Lemma~\ref{lem:block-outer-first-moment-interface}, fix a prefix cutoff \(H\), an episode cutoff
\(R\), and a live-slot cutoff \(\theta N\).  Suppose the prefix satisfies
\[
P_H:=\sum_{1\le h\le H} A_h^{\mathrm{out}}p_h^{\mathrm{in}}(\delta)\le P,
\]
and suppose the residual \(h>H\) is bounded by four audited pieces
\[
\begin{array}{ll}
T_{\mathrm{lo},\le R} & (1\le r\le R,\ M\le\theta N),\\
T_{\mathrm{hi},\le R} & (1\le r\le R,\ M>\theta N),\\
T_{\mathrm{lo},>R} & (r>R,\ M\le\theta N),\\
T_{\mathrm{hi},>R} & (r>R,\ M>\theta N).
\end{array}
\]
Then
\[
\Pr[d_{\min}(\Csc)\le \lfloor\delta N\rfloor]
\le
P+T_{\mathrm{lo},\le R}+T_{\mathrm{hi},\le R}+T_{\mathrm{lo},>R}+T_{\mathrm{hi},>R}.
\]
In particular, a finite certificate proving that the right-hand side is at most \(2^{-\lambda}\) gives failure
probability at most \(2^{-\lambda}\).
\end{lemma}

\begin{proof}
This is only bookkeeping.  Lemma~\ref{lem:block-outer-first-moment-interface} reduces the failure probability to the
first moment over nonzero outer weights.  The prefix contributes \(P_H\), and the residual partition covers all
post-prefix episode configurations.  Summing the audited upper bounds gives the displayed estimate.
\end{proof}

Thus a local code with a published or certified spectrum drops into exactly the same first-moment interface
\(\sum_h A_h^{\mathrm{out}}p_h^{\mathrm{in}}(\delta)\).  The asymptotic theorem still uses the random banded outer; the
block code is best viewed as a finite-\(n\) implementation variant whose proof burden is shifted to a short local
spectrum certificate.


\subsection{Full-split local-code recursive inner}
\label{subsec:fullsplit-local-code-inner}

The finite certificate below instantiates the inner interface with the full-split extended-BCH block inner used in the
current checkpoint.  Its local-code-independent input is the short block spectrum and the random split law; the concrete
checked theorem uses the EBCH \([128,64,22]\) spectrum.

\paragraph{Full-codeword random-split variant.}
There is a proof-friendlier variant that avoids the systematic split profile \(B_{j,q'}\).  After forming
\[
V_i=U_i+S_{i-1},
\]
if \(V_i\ne0\), apply a block scrambler \(A_i\in GL_b(2)\), encode \(Z_i=A_iV_i\) with the rate-half BCH block code,
and randomly split the \(2b\) BCH coordinates into \(b\) emitted coordinates and \(b\) state coordinates:
\[
C_i=E_{\rm BCH}(A_iV_i),\qquad (Y_i,S_i)=\operatorname{split}_b(\tau_i C_i).
\]
For random \(A_i\), conditional on \(V_i\ne0\), the BCH codeword is uniform over nonzero BCH codewords.  Hence the
ordinary full BCH weight spectrum \(A_w\) is sufficient.  If \(w=\wt(C_i)\), then
\[
\Pr[\wt(S_i)=q'\mid w]
=
\frac{\binom{b}{q'}\binom{b}{w-q'}}{\binom{2b}{w}},
\qquad
\wt(Y_i)=w-q'.
\]
The user's ``self turnoff'' concern is exactly the \(q'=0\) tail of this hypergeometric split:
\[
\Pr[q'=0\mid w]=\frac{\binom{b}{w}}{\binom{2b}{w}},
\]
and is therefore directly charged rather than hidden in a split-profile theorem.  Using the published
EBCH \([128,64]\) spectrum table, the unconditional nonzero-codeword split tails are
\[
\begin{array}{c|c}
\text{tail} & \log_2\Pr[\text{tail}]\\
\hline
q'=0 & -64.004574\\
q'\le4 & -44.626752\\
q'\le8 & -31.743503\\
q'\le16 & -14.658598
\end{array}
\]
Thus exact self-termination is tiny, while small-state tails remain explicit terms in the weight-chain transfer.
There is also a useful exact local simplification.  Suppose the incoming state has support size \(q\), and the next
input block contributes \(r\) permuted nonzero coordinates.  The only way to turn off before applying the BCH branch is
exact cancellation of the two subsets.  Hence
\[
\Pr[V_i=0\mid q,r]
=
\begin{cases}
\binom{b}{q}^{-1}, & q=r,\\
0, & q\ne r.
\end{cases}
\]
On the complementary event \(V_i\ne0\), the next \((\wt(Y_i),\wt(S_i))\) law is the same fixed full-spectrum split law,
independent of \(q\), \(r\), and \(\wt(V_i)\).  Thus the transfer kernel is not a general two-parameter local profile:
it is an exact-cancellation atom plus one weighted BCH split vector.  This is the main reason the full-split variant is
cleaner than the systematic split model.

The first bucketed diagnostics for this full-split variant are encouraging.  Against the same RM outer prefix at
\(N=2^{21}\), \(\delta=.09\), \(L=5949\), and the EBCH full spectrum, the middle bucket
\(4001\text{--}8000\) with \(r=1\), \(16\)-weight state buckets, gives
\[
\begin{array}{c|c}
h & \log_2\text{ term}\\
\hline
32 & -23.40\\
64 & -30.64\\
80 & -27.80
\end{array}
\]
using \(\lambda=0.0002\) for the latter two rows.  The \(h=32\) full coarse range \(1\text{--}12000\) is still
dominated by the near-boundary placement bucket and gives about \(-17.38\) bits, so the boundary split remains real.
However, unlike the systematic-split diagnostic, the intermediate \(h\) rows do not appear to drift toward failure in
the tested bucket.

After exploiting the exact-cancellation-plus-fixed-branch structure, the diagnostic can run with exact state weights
\(q=0,\ldots,b\) rather than \(16\)-weight buckets.  The exact-\(q\) ridge check for
\(h\in\{32,40,64,80\}\), \(r=1\), and the three coarse gap buckets
\(1\text{--}4000\), \(4001\text{--}8000\), \(8001\text{--}12000\) gives total exponent about \(-17.56\) bits over
this table.  The largest rows are
\[
\begin{array}{c|c|c|c}
h & r & \text{gap bucket} & \log_2\text{ term}\\
\hline
32 & 1 & 1\text{--}4000 & -17.88\\
40 & 1 & 1\text{--}4000 & -19.91\\
64 & 1 & 1\text{--}4000 & -26.21\\
80 & 1 & 4001\text{--}8000 & -28.33\\
80 & 1 & 1\text{--}4000 & -30.41
\end{array}
\]
For \(h=32\), exact-\(q\) and \(r=1,2,3,4\) over the same three buckets gives total exponent about \(-17.88\) bits,
again dominated by the \(r=1\), \(1\text{--}4000\) boundary bucket.  In that near-boundary bucket the \(r=2,3,4\)
terms are approximately \(-27.18\), \(-37.14\), and \(-47.58\), respectively.  Thus the present finite evidence says
that the full-split bottleneck has moved from intermediate \(h\) to the first nonzero RM weight, one-coordinate
first-active, boundary-placement branch.

This is still a finite diagnostic, not a proof.  The proof target should now be an exact-\(q\) first-moment transfer
using the cancellation formula above, followed by a tail domination lemma showing that the remaining \(h\)'s, larger
first-block occupancies, and farther gap buckets are dominated by the checked ridge.  The diagnostic script is
\[
\texttt{python scripts/certify\_fullsplit\_inner\_bucket\_sum.py}.
\]

A subsequent prefix-summation check shows that the exact-\(q\) transfer cannot simply be used as one unconditional
black-box exponent.  If one reuses the measured bucket exponents
\[
-0.475,\qquad -37.407,\qquad -74.323
\]
for the three buckets \(1\text{--}4000\), \(4001\text{--}8000\), and \(8001\text{--}12000\), and sums the RM outer
prefix through \(h=500\), the far bucket is not summable; the pessimistic prefix sum is strongly positive and peaks at
the artificial endpoint \(h=500\).  This should not be read as a counterexample to the construction, but it is a real
warning about the proof method.  The one-shot Chernoff moment is capped by the \(q'=0\) self-turnoff atom: even a long
live trajectory has a rare early-termination branch, and the unconditional exponential moment pays for that branch
poorly when summed against many outer words.  Thus the full proof should split the inner event into survival/episode
classes.  On a long-survival class the emitted BCH split weights should give a much stronger Chernoff exponent; on a
short-episode class the proof must charge the probability of self-turnoff and then restart with the remaining input
placement.  In other words, the full-split variant removes the unknown systematic split profile, but it does not remove
the need for an episode-count or survival-conditioned correction.  The current bookkeeping helper is
\[
\texttt{python scripts/sum\_fullsplit\_uniform\_inner.py}.
\]

The survival-conditioned side of this split has very large slack.  Let
\[
p_0:=\Pr[q'=0]
\]
for one nonzero BCH split branch.  For the EBCH \([128,64]\) table,
\[
\log_2 p_0=-64.004574.
\]
If the \(q'=0\) atom is excluded for \(T\) consecutive live branches, the one-branch split law has essentially the same
mean output \(32\), but no early absorbing off state.  Optimizing the Chernoff parameter for the three current live
lengths gives
\[
\begin{array}{c|c|c}
T & \log_2\text{ survival-conditioned Chernoff bound} & \lambda\\
\hline
5949 & -20.01 & .017\\
9949 & -78398.76 & .864\\
13949 & -228370.90 & 1.316
\end{array}
\]
Using these as bucket exponents and summing the RM prefix through \(h=500\) gives total about \(-34.77\) bits, dominated
by the near-boundary bucket.  Thus the long-survival branch is not the problem; the middle and far buckets essentially
disappear once self-turnoff is excluded.

The short-episode branch still needs real work.  A crude union charge \(T p_0\) for one self-turnoff, followed by
discarding all remaining output, is far too weak: the same \(h\le500\) prefix remains strongly positive and again peaks
at the artificial \(h=500\) endpoint.  The missing lemma must therefore be recursive or episode-counting: if an early
episode self-terminates, the proof must restart with the remaining input placement and charge either another
self-turnoff, a long-survival Chernoff event, or the event that all remaining input has moved into the late window.  The
supporting diagnostic for the survival side is
\[
\texttt{python scripts/analyze\_fullsplit\_survival.py}.
\]

There is a clean candidate for this episode-counting lemma.  Fix the first active block and let \(T\) blocks remain
after it.  Suppose \(H\) input symbols remain.  If exactly \(x\) of the \(T\) later blocks are occupied by these
symbols, then the \(x\) occupied blocks and the terminal suffix define \(m=x+1\) zero-gaps after active blocks,
including the gap after the first active block.  The occupancy law is exact:
\[
\Pr[X=x]
=
\frac{\binom{T}{x}\,[z^H]\big((1+z)^b-1\big)^x}{\binom{bT}{H}}.
\]
Conditional on \(X=x\), the gap vector is a uniform composition of \(G=T-x\) empty blocks into \(m=x+1\) nonnegative
parts.  Therefore the expected number of \(e\)-subsets of gaps whose total skipped length is \(s\) is
\[
\binom{m}{e}\,
\frac{\binom{s+e-1}{e-1}\binom{G-s+m-e-1}{m-e-1}}
{\binom{G+m-1}{m-1}},
\]
with the obvious endpoint conventions for \(e=0\) and \(e=m\).  If those \(e\) gaps are skipped by \(q'=0\)
self-turnoffs, their probability contribution is \(p_0^e\).  The remaining \(T-s\) live branches are then charged by
the survival-conditioned Chernoff factor
\[
\inf_{\lambda>0}\exp(\lambda d)
\left(\sum_{j,q>0}p_{j,q}e^{-\lambda j}\right)^{T-s},
\]
where \(p_{j,q}\) is the full BCH split law.  This gives a completely explicit self-turnoff episode upper bound:
\[
\sum_x \Pr[X=x]\sum_{e,s}
p_0^e\,
\mathbb{E}\big[\#\{e\text{-gap subsets with skipped length }s\}\mid X=x\big]\,
\inf_{\lambda>0}e^{\lambda d}M_+(\lambda)^{T-s}.
\]
It is still missing exact input/state cancellation events \(V_i=0\), but it directly addresses the dominant
self-turnoff/restart mechanism that broke the one-shot MGF.

The first diagnostics for this explicit episode-gap bound are favorable.  They use \(r=1\) in the first active block,
so \(H=h-1\), and charge only self-turnoffs \(q'=0\).  The near bucket \(T=5949\) is dominated by the no-self-turnoff
survival term and gives about \(-20.01\) inner bits for both \(h=32\) and \(h=80\).  For the middle and far buckets,
the bound strengthens with \(h\):
\[
\begin{array}{c|c|c}
h & T=9949 & T=13949\\
\hline
32 & -82.30 & -97.42\\
80 & -117.05 & -155.57\\
160 & -176.17 & -253.70\\
320 & -295.41 & -450.97\\
500 & -430.02 & -673.38
\end{array}
\]
The dominant middle/far terms in these samples select a single large skipped gap (\(e=1\)); allowing up to four skipped
gaps did not change the \(h=500\) middle or far bound.  This is the missing shape: high \(h\) supplies many occupied
blocks, so the probability of a very large skipped zero-gap decays rapidly enough to overcome the outer growth.
The diagnostic is
\[
\texttt{python scripts/bound\_fullsplit\_episode\_gaps.py}.
\]

The remaining exact-cancellation hazard appears smaller than the self-turnoff atom in the finite range currently being
certified.  If a live branch leaves state weight \(q\), then at the next block
\[
\Pr[V_i=0\mid q,H,T]
=
\frac{\binom{b(T-1)}{H-q}}{\binom{bT}{H}},
\]
where \(H\) input symbols remain over \(T\) blocks.  Averaging this over the full BCH split state-weight law gives the
following one-step exact-hit exponents:
\[
\begin{array}{c|c|c|c}
H & T=5949 & T=9949 & T=13949\\
\hline
31 & -71.59 & -72.33 & -72.82\\
79 & -70.25 & -70.98 & -71.47\\
159 & -69.25 & -69.98 & -70.46\\
319 & -68.26 & -68.99 & -69.47\\
499 & -67.64 & -68.35 & -68.83
\end{array}
\]
The worst displayed value occurs at the smallest remaining span.  More generally, for fixed \(H,q\), the ratio
\(\binom{b(T-1)}{H-q}/\binom{bT}{H}\) is decreasing in \(T\): adding more outside coordinates while keeping \(H\)
fixed only lowers the chance that the distinguished next block contains exactly the prescribed \(q\)-set and no other
selected coordinate.  Thus for the finite prefix \(H\le499\) and \(T\ge5949\), scanning the left endpoint gives
\[
\log_2 p_0=-64.00457432724919,\qquad
\log_2\eta=-67.63641076470455,
\]
and hence
\[
\log_2(p_0+\eta)=-63.89264922047382.
\]
The finite certificate uses the slightly rounded-up atom
\[
\log_2 p_{\rm term}\le -63.8926492.
\]
Rerunning the \(h=500,T=13949\) episode-gap diagnostic with this inflated atom changes the inner exponent from
\(-673.38\) to \(-673.27\), and the near \(h=32,T=5949\) ridge is unchanged because it is dominated by the
no-termination survival branch.  Thus exact input/state cancellation can be folded into the same skipped-gap ledger,
at least for the current \(h\le500\) finite prefix, by replacing \(p_0\) with
\[
p_{\mathrm{term}}:=p_0+\eta,\qquad
\eta=\sup_{H,T}\sum_{q>0}\Pr[Q=q]\frac{\binom{b(T-1)}{H-q}}{\binom{bT}{H}}.
\]
The audit-grade cancellation verifier is
\[
\texttt{python scripts/verify\_fullsplit\_turnoff\_atom.py}.
\]

For theorem-facing wrappers outside this finite prefix we use a slightly looser global atom.  Let \(n=bT\),
\(\alpha=H/n\), and \(T\ge5949\).  For a fixed \(q\)-set in the next block,
\[
\frac{\binom{n-b}{H-q}}{\binom{n}{H}}
=\frac{(H)_q(n-H)_{b-q}}{(n)_b}
\le
\left(1-\frac{b-1}{b\cdot5949}\right)^{-b}
\alpha^q(1-\alpha)^{b-q}.
\]
Averaging over \(Q\) gives a Bernstein polynomial:
\[
\sum_{q>0}\Pr[Q=q]\alpha^q(1-\alpha)^{b-q}
=
\sum_{q>0}\frac{\Pr[Q=q]}{\binom bq}B_{q,b}(\alpha).
\]
For the EBCH split law the largest nonzero Bernstein coefficient is the \(q=b\) coefficient, giving
\[
\log_2\eta_{\rm glob}\le -62.98700592417877,\qquad
\log_2(p_0+\eta_{\rm glob})\le -62.4078758.
\]
The verifier is
\[
\texttt{python scripts/verify\_fullsplit\_global\_turnoff\_envelope.py}.
\]

\begin{lemma}[Certified full-split termination atoms]
\label{lem:fullsplit-certified-termination-atoms}
For the EBCH \([128,64]\) full-split law with \(b=64\), let \(Q\) denote the next-state weight after one nonzero BCH
split branch.  The self-turnoff atom satisfies
\[
\log_2\Pr[Q=0]=-64.00457432724919.
\]
For every finite-prefix ledger state with \(0\le H\le499\) remaining input ones and \(T\ge5949\) remaining blocks, the
episode-ledger exposure described in Lemma~\ref{lem:fullsplit-selected-gap-product} has a valid charged termination atom
\[
\log_2 p_{\rm term}^{\rm fin}\le -63.8926492.
\]
For every \(0\le H\le bT\) and \(T\ge5949\), the stronger global exposure has a valid charged termination atom
\[
\log_2 p_{\rm term}^{\rm glob}\le -62.4078758.
\]
The global atom is also safe after conditioning on any fixed next-block target support.
\end{lemma}

\begin{proof}
The \(Q=0\) value is read from the EBCH split law.  For exact input/state cancellation in the finite-prefix ledger, a
fixed state support of size \(q\) is hit by the next distinguished block with probability
\[
\frac{\binom{b(T-1)}{H-q}}{\binom{bT}{H}},
\]
and this ratio is decreasing in \(T\) for fixed \(H,q\).  Averaging over \(Q>0\), scanning \(0\le H\le499\) at
\(T=5949\) gives
\[
\log_2\eta_{\rm fin}\le -67.63641076470455.
\]
Thus
\[
\log_2\bigl(\Pr[Q=0]+\eta_{\rm fin}\bigr)
=-63.89264922047382
\le -63.8926492.
\]

For the global atom, write \(n=bT\) and \(\alpha=H/n\).  For a fixed \(q\)-set in the next block,
\[
\frac{\binom{n-b}{H-q}}{\binom{n}{H}}
=\frac{(H)_q(n-H)_{b-q}}{(n)_b}
\le
\left(1-\frac{b-1}{b\cdot5949}\right)^{-b}\alpha^q(1-\alpha)^{b-q}.
\]
Averaging over \(Q>0\) gives a Bernstein polynomial.  The verifier checks the finite coefficient inequality
\[
\max_{1\le q\le b}\frac{\Pr[Q=q]}{\binom bq}=\Pr[Q=b],
\]
so the exact-cancellation envelope is
\[
\log_2\eta_{\rm glob}\le -62.98700592417877.
\]
Adding the \(Q=0\) self-turnoff atom gives
\[
\log_2\bigl(\Pr[Q=0]+\eta_{\rm glob}\bigr)
=-62.407875819071386
\le -62.4078758.
\]
The same coefficient inequality also proves the fixed-support claim, since a fixed target support of size \(q\) is hit
with probability \(\Pr[Q=q]/\binom bq\), which is at most the displayed maximum.
\end{proof}

\paragraph{First-active placement coverage.}
For the finite checkpoint \(N=2^{21}\), \(b=64\), write \(M=N/b=32768\).  If the first active inner block is
\(i_1\), define
\[
T:=M-i_1,\qquad r:=\wt(U_{i_1}),\qquad H:=h-r.
\]
Thus \(T\) is the number of later \(b\)-blocks after the first active block.  The current checked ledger uses the split
\[
\begin{array}{c|c|c}
\text{regime} & T\text{ range} & \text{current status}\\
\hline
\text{ultra-late} & 0\le T<5949 & \text{placement cap checked}\\
\text{audited window} & 5949\le T\le17948 & \text{episode ledger checked}\\
\text{early} & T>17948 & \text{episode ledger checked}
\end{array}
\]
The audited window is exactly the script parameterization
\[
T=5949+g-1,\qquad 1\le g\le12000,
\]
split into gap buckets \(1\text{--}4000\), \(4001\text{--}8000\), and \(8001\text{--}12000\).  Since \(r\le64\), the
far endpoint \(T=17948\) implies the feasibility cutoff
\[
h=H+r\le 64\cdot17948+64=1148736.
\]
The verifier
\[
\texttt{python scripts/verify\_fullsplit\_finite\_ledger.py}
\]
now also checks the placement-only ultra-late prefix
\[
\sum_{32\le h\le500} A_h^{\rm RM}
\Pr[\text{first active block lies in the final }5949\text{ blocks}]
\le 2^{-41.113442},
\]
using exact RM direct-sum outer coefficients through \(h=500\).  For \(h\ge501\), the same placement cap is handled by
fixed-pole interval rows in Corollary~\ref{cor:fullsplit-postprefix-checked-ledger}; these rows contribute at most
\(-545.690526\) bits.  The window ledger itself is
\[
\log_2\mu_{\rm window}^{\rm exact\ outer}\le -37.383345,
\]
and the union of this checked window with the checked ultra-late prefix has
\[
\log_2\mu_{\rm ultra\mbox{-}late\ prefix\ +\ window}^{\rm exact\ outer}\le -37.278528.
\]
The early \(h\le500\) part is now checked separately using the gap buckets
\[
12001\text{--}17000,\qquad 17001\text{--}22000,\qquad 22001\text{--}26819,
\]
equivalently \(17949\le T\le32767\).  The explicit \(e\le16\) bucket-uniform survival ledger gives
\[
\log_2\mu_{32\le h\le500,\ {\rm early}}^{e\le16,\ {\rm exact\ outer}}\le -86.910456,
\]
and the crude \(e\ge17\) remainder gives
\[
\log_2\mu_{32\le h\le500,\ {\rm early}}^{e\ge17,\ {\rm exact\ outer}}\le -319.977808.
\]
Thus the combined \(h\le500\) all-first-active-position value remains
\[
\log_2\mu_{32\le h\le500}^{\rm all\ first\ active\ positions,\ exact\ outer}\le -37.278528,
\]
still dominated by the audited window and ultra-late prefix.  The proof debt is therefore no longer the boundary slice
\(T\approx5949\), nor the small-weight early slice.  The checked early rows now cover
\[
501\le h\le N/2,\qquad T>17948,
\]
with total \(-380.829185\) bits to displayed precision.  This combines the original endpoint table through
\(h=75000\), an accelerated endpoint recurrence through \(h=350000\), and a paired-\(T\) high-density recurrence
through \(h=N/2\).  The post-prefix part of the ultra-late cap beyond \(h=500\) is separately covered by the
placement-only interval rows below.

Putting the inflated termination atom into the episode-gap sweep and multiplying by the RM outer and exact first-active
placement gives the first combined \(h\)-dependent checkpoint.  For \(r=1\), gap buckets
\(1\text{--}4000\), \(4001\text{--}8000\), \(8001\text{--}12000\), and
\[
h\in\{32,80,160,320,500\},
\]
the total over these sampled rows is about \(-37.41\) bits and is dominated by the near-boundary \(h=32\) row:
\[
\begin{array}{c|c|c|c}
h & \text{gap bucket} & \log_2 p_{\rm inner}^{\rm episode} & \log_2\text{ full term}\\
\hline
32 & 1\text{--}4000 & -20.01 & -37.41\\
80 & 1\text{--}4000 & -20.01 & -49.94\\
160 & 1\text{--}4000 & -20.01 & -73.44\\
32 & 4001\text{--}8000 & -82.19 & -83.99\\
32 & 8001\text{--}12000 & -97.30 & -87.46\\
80 & 4001\text{--}8000 & -116.93 & -107.86\\
80 & 8001\text{--}12000 & -155.45 & -117.27\\
320 & 1\text{--}4000 & -20.01 & -134.87\\
500 & 1\text{--}4000 & -20.01 & -204.02
\end{array}
\]
The omitted sampled middle/far \(h=160,320,500\) rows are lower still.  This is not yet a proof over all \(h\) and \(r\),
but it confirms the intended decomposition: the near bucket is controlled mostly by placement plus the survival
Chernoff at \(T=5949\), while the middle/far buckets require the episode-gap skipped-gap count and then become more
negative as \(h\) grows.  The combined sweep wrapper is
\[
\texttt{python scripts/sweep\_fullsplit\_episode\_terms.py}.
\]

A faster envelope summation gives the first dense-prefix checkpoint.  Using the sampled episode bounds above as a
piecewise-linear working envelope in \(H=h-r\), and summing exactly over
\[
32\le h\le 500,\qquad 1\le r\le 64,
\]
for the same three gap buckets gives
\[
\log_2 \sum_{\substack{32\le h\le500\\1\le r\le64\\\text{three buckets}}}
A_h^{\rm RM}\,
\Pr[\text{first active occupancy }r\text{ in bucket}]\,
p_{\rm inner}^{\rm env}(h-r,\text{bucket})
=-34.767174.
\]
The bucket contributions are
\[
\begin{array}{c|c}
\text{gap bucket} & \log_2\text{ contribution}\\
\hline
1\text{--}4000 & -34.767174\\
4001\text{--}8000 & -82.122759\\
8001\text{--}12000 & -85.867481
\end{array}
\]
and the first occupancy split begins
\[
\begin{array}{c|c}
r & \log_2\text{ contribution}\\
\hline
1 & -34.769819\\
2 & -43.861905\\
3 & -53.569247\\
4 & -63.713307
\end{array}
\]
In the smoothed Cauchy-outer ledger, the peak row is still \(h=32,r=1\), gap \(1\text{--}4000\), with inner exponent
\(-20.011735\) and full term \(-37.385767\).  Before using the exact RM support, the full prefix is not a shapeless
table effect.  It is a narrow near-boundary ridge:
\[
\begin{array}{c|c}
\text{slice} & \text{share of the }e\le8,\ 32\le h\le500\text{ prefix}\\
\hline
\text{gap }1\text{--}4000 & 1-O(10^{-14})\\
r=1 & .998168813\\
h\le64 & .997505683\\
h\le80 & .999863886\\
h\le96 & .999992575
\end{array}
\]
The smoothed total prefix value is \(2.618593\) bits above the single peak row.  Along the dominant \(r=1\), first-gap ridge,
the largest adjacent ratio is \(0.8526993\) from \(h=32\) to \(33\), and after that it is at most \(0.8337949\) in the
audited range.  The smoothed finite version already has a summable form: replacing the ridge by the infinite
geometric bound
\[
T_{32}\left(1+\frac{0.852700}{1-0.833795}\right)
\]
and adding the exact non-ridge remainder gives
\[
\log_2\mu_{\rm prefix}^{\rm ratio}\le -34.767141,
\]
only \(3.4\cdot10^{-5}\) bits above the exact finite prefix ledger.  The decomposition helper is
\[
\texttt{python scripts/analyze\_fullsplit\_prefix\_ridge.py}.
\]
The ratio itself factors into outer-spectrum growth, first-active placement, and the inner \(e\le8\) value.  On this
ridge the inner span is only \(3.2\cdot10^{-12}\) bits.  If
\[
S_H:=\sum_{g=1}^{4000}\binom{64(5949+g-1)}{H},
\]
then the \(r=1\), first-gap placement contribution satisfies
\[
\frac{P_{h+1}}{P_h}
=
\frac{S_h}{S_{h-1}}\frac{h+1}{N-h}.
\]
The current finite component certificate checks
\[
\begin{array}{c|ccc}
\text{step} & A_{h+1}^{\rm out}/A_h^{\rm out} & P_{h+1}/P_h & \text{inner ratio}\\
\hline
32\to33 & 2.808686 & .303594 & \le 1.000000001\\
h\ge33 & 2.746420 & .303594 & \le 1.000000001.
\end{array}
\]
The placement entry is now checked independently of the ledger CSV by an exact-integer cross-multiplication certificate:
\[
\max_{32\le h<500}\frac{S_h}{S_{h-1}}\frac{h+1}{N-h}
\le .303594,
\]
where \(.303594=151797/500000\).  The certificate performs \(468\) exact comparisons, finds the maximum at
\(h=32\), and verifies a positive rational slack of about \(2^{-21.8672}\) below the threshold.  The exact peak is
\(0.303593738594\).  The naive endpoint bound
\[
\frac{\binom{64(5949+4000-1)}{h}}{\binom{64(5949+4000-1)}{h-1}}\frac{h+1}{N-h}
\]
is \(0.3130655524\) at \(h=32\), about \(1.0312\) times too large, so the proof must use the averaged-binomial sum
\(S_H\) rather than endpoint domination if one works with the smoothed outer envelope.

The exact RM direct-sum support gives a sharper and cleaner finite certificate.

\begin{lemma}[Exact RM prefix support split certificate]
\label{lem:rm-prefix-support-split}
Consider the finite checkpoint \(N=2^{21}\), \(\delta=.09\), and the RM direct-sum outer consisting of \(4096\)
copies of the local \([512,256]\) RM block whose weight enumerator is
\[
W_{\rm loc}(z)=\sum_{a=0}^{512} A_a^{\rm loc}z^a.
\]
Let
\[
A_h^{\rm RM}:=[z^h]\,W_{\rm loc}(z)^{4096}.
\]
The integer-truncated coefficient certificate through \(h\le500\) verifies that the positive support begins
\[
32,\ 48,\ 56,\ 60,\ 64,\ldots,
\]
has no positive weight below \(32\), has no positive weight in \(33\le h\le47\), and has \(114\) positive supported
weights in \(1\le h\le500\).  The coefficient build uses exact truncated polynomial exponentiation: for the present
\(4096=2^{12}\) direct sum, the certificate records one multiply step and twelve squaring steps.  The \(h=32\)
coefficient has \(38\) bits, while the next supported coefficient at \(h=48\) has \(52\) bits.  Replacing the smoothed
outer rows in the audited
\(32\le h\le500\) ledgers by the exact coefficients \(A_h^{\rm RM}\) gives
\[
\begin{array}{c|ccc}
\text{piece} & \log_2\text{ total} & \log_2(h=32) & \log_2(h>32)\\
\hline
\text{window }5949\le T\le17948 & -37.383345 & -37.383482 & -50.738253\\
\text{ultra-late }T<5949 & -41.113442 & -41.113442 & -66.416502
\end{array}
\]
The same exact-outer reweighting gives
\[
\log_2\mu_{32\le h\le500,\ {\rm early},\ e\le16}^{\rm exact\ outer}\le -86.910456,
\qquad
\log_2\mu_{32\le h\le500,\ {\rm early},\ e\ge17}^{\rm exact\ outer}\le -319.977808.
\]
\end{lemma}

\begin{proof}
The local enumerator coefficients \(A_a^{\rm loc}\) are read as integers from the checked CSV spectrum file.  Since the
outer is a direct sum of \(4096\) identical blocks, its weight enumerator is \(W_{\rm loc}(z)^{4096}\).  The certificate
computes this polynomial modulo \(z^{501}\) by binary exponentiation with exact integer convolutions, so no
floating-point or smoothing step enters the support calculation.  The split values are then obtained by reweighting each
audited CSV row by
\[
\log_2 A_h^{\rm RM}-\log_2 A_h^{\rm smooth}
\]
and then partitioning the rows into the \(h=32\) slice and the supported \(h>32\) remainder before log-summing.  The
verifier
\[
\texttt{python scripts/verify\_fullsplit\_finite\_ledger.py}
\]
checks the first two positive support weights \(32,48\), the displayed split totals, and the final combined totals.
\end{proof}

Thus the apparent adjacent \(h=33,34,\ldots\) ridge is a Cauchy-smoothing artifact.  Reweighting the same audited
inner/placement rows by exact RM direct-sum coefficients gives
\[
\log_2\mu_{32\le h\le500,\ {\rm window}}^{\rm exact\ outer}\le -37.383345,
\]
with the \(h=32\) slice alone contributing \(-37.383482\) and all supported \(h>32\) rows together contributing only
\(-50.738253\).  The certificate also checks that there are \(114\) positive supported weights in this prefix and that
the next supported weight after \(32\) is \(48\).
Within the \(h=32\) slice, the dominant row is exactly
\[
h=32,\qquad r=1,\qquad 1\le g\le4000,\qquad T=5949,
\]
with exact-outer term \(-37.385767\).  The full exact-outer window prefix is only \(0.002422\) bits above this
single row.  Removing this row leaves total log contribution \(-46.602718\); the remaining \(h=32\) part is
\(-46.687229\), while all supported \(h>32\) rows are \(-50.738253\), i.e. \(13.352486\) bits below the dominant row.
The exact ultra-late prefix is
\[
\log_2\mu_{32\le h\le500,\ T<5949}^{\rm exact\ outer}\le -41.113442,
\]
with supported \(h>32\) contributing only \(-66.416502\).
The combined finite ledger becomes
\[
\log_2\mu_{\rm finite}^{\rm exact\ outer}\le -37.383345,\qquad
\log_2\mu_{\rm late+window}^{\rm exact\ outer}\le -37.278528.
\]
Thus the current theorem target is no longer an analytic adjacent-\(h\) ridge for the smoothed outer bound; it is the
exact-support \(h=32\) term plus a support-weight remainder.
The standard verifier
\[
\texttt{python scripts/verify\_fullsplit\_finite\_ledger.py}
\]
checks both the smoothed ratio skeleton and the exact-support replacement.  The exact RM prefix helper is
\[
\texttt{python scripts/certify\_rm\_outer\_prefix\_exact.py},
\]
and the placement subcheck can be run alone as
\[
\texttt{python scripts/certify\_prefix\_placement\_ratio.py}.
\]
This is not yet a theorem, because the interpolation envelope between sampled \(H\)-knots must be justified or replaced
by a grid-plus-monotonicity certificate.  It does show that the prefix no longer has a broad hidden intermediate-\(h\)
ridge once the episode-gap correction is included.  The envelope summation helper is
\[
\texttt{python scripts/sum\_fullsplit\_episode\_envelope.py}.
\]

The dominant middle/far episode term has \(e=1\) selected skipped gap, and it admits the closed form
\[
\Pr[X=x]\,
\mathbb{E}[\#\{\text{one skipped gap of length }s\}\mid X=x]
=
\frac{[z^H]((1+z)^b-1)^x}{\binom{bT}{H}}\,(x+1)\binom{T-s-1}{x-1}.
\]
This fast formula matches the previous \(e\le4\) knot values at the tested points, because the \(e=0\) or \(e=1\) term
dominates there.  The implementation now factors the skipped-gap suffix sum, which depends on \(x\) and \(T\) but not
on \(H\), so denser interpolation audits are cheap.  It also includes the \(H=0\) endpoint explicitly: if no later input
ones remain, the \(e=1\) branch is the single full zero gap and contributes exactly the effective termination atom
\[
\log_2 p_{\rm term}\le -63.8926492.
\]
Using full integer \(H=0,\ldots,499\) CSVs, with extrapolation disallowed in the envelope summation, removes the
interpolation issue from the finite prefix.  The smoothed-outer \(e\le8\) ledger gives the total prefix value
\[
-34.767174,
\]
with middle/far bucket contributions
\[
4001\text{--}8000:\ -82.122759,\qquad
8001\text{--}12000:\ -85.867481.
\]
Reweighting the same rows by exact RM direct-sum outer coefficients gives the current finite value
\[
\log_2\mu_{32\le h\le500,\ e\le8}^{\rm exact\ outer}\le -37.383345.
\]
\begin{lemma}[Tiny-prefix interior-\(T\) finite certificate]
\label{lem:fullsplit-tiny-prefix-interiorT}
Consider the finite checkpoint \(N=2^{21}\), \(\delta=.09\), \(b=64\), and the finite-prefix termination atom
\(\log_2 p_{\rm term}\le -63.8926492\).  Let
\[
{\cal B}_1=[5949,9948],\qquad
{\cal B}_2=[9949,13948],\qquad
{\cal B}_3=[13949,17948]
\]
be the three residual-length buckets corresponding to gap ranges \(1\text{--}4000\),
\(4001\text{--}8000\), and \(8001\text{--}12000\).  For \(0\le H\le499\), let
\(p_{\le8}(H,T)\) denote the exact full-split selected-gap inner probability with at most eight charged
terminations, computed from the closed-form fixed-\(e\) episode summands.  The tracked finite audit certifies that,
for each bucket \({\cal B}_i=[T_i^-,T_i^+]\),
\[
\log_2 p_{\le8}(H,T)
\le
\log_2 p_{\le8}(H,T_i^-)+10^{-7}
\qquad
(0\le H\le499,\ T_i^-\le T\le T_i^+).
\]
The recorded maximum excess is \(2.19205276153\cdot 10^{-9}\), attained at the left endpoint itself; the largest
left-endpoint reproduction error is \(5.52478240934\cdot 10^{-9}\) bits.  Consequently, for the current finite
first-moment ledger, the left-endpoint CSVs at \(T=5949,9949,13949\) are valid representatives for the whole
\(32\le h\le500,\ e\le8\) prefix up to the stated numerical tolerance.
\end{lemma}

\begin{proof}
For fixed \(H\), \(T\), and selected-gap count \(e\), the exact episode summand is finite.  The skipped-gap suffix
\[
\sum_{L=x}^{T}\binom{T-L+e-1}{e-1}\binom{L-e}{x-e}\,S(L)
\]
depends on \(x,e,T\), but not on \(H\).  The audit implementation is
\[
\begin{gathered}
\texttt{scripts/check\_fullsplit\_exact\_grid\_interior.py}\\
\texttt{--h-values 0:499,\quad --sample-step 1,\quad --method combined}.
\end{gathered}
\]
It sweeps all interior \(T\)'s by maintaining the short cumulative sums for these suffixes, rather than regenerating a
separate CSV for every \(T\).  It checks all \(500\) values of \(H\) in all three buckets, writing one summary row for
each \((H,{\cal B}_i)\), hence \(1500\) rows.  The standard ledger verifier
\[
\texttt{python scripts/verify\_fullsplit\_finite\_ledger.py}
\]
now reads the tracked artifact
\[
\texttt{scripts/fullsplit\_exact\_interior\_h0\_499\_allT.csv}
\]
and fails unless the row count is \(1500\), the worst interior excess is at most \(10^{-7}\), and the worst
left-endpoint reproduction error is at most \(5\cdot10^{-6}\).  It also checks that the artifact is exactly the
rectangle \(H=0,\ldots,499\) crossed with the three displayed buckets, that each row covers all \(4000\) integer
values of \(T\) in its bucket, and that every row is maximized at the bucket's left endpoint.  The current artifact has
the numerical maxima stated above, leaving \(9.78079472385\cdot10^{-8}\) bits of excess slack and
\(4.99447521759\cdot10^{-6}\) bits of reproduction-error slack against the verifier thresholds.
\end{proof}

The same closed form extends to any fixed number \(e\) of selected skipped gaps:
\[
\Pr[X=x]\,
\mathbb{E}[\#\{e\text{ skipped gaps of total length }s\}\mid X=x]
=
\frac{[z^H]((1+z)^b-1)^x}{\binom{bT}{H}}\,
\binom{x+1}{e}\binom{s+e-1}{e-1}\binom{T-s-e}{x-e},
\]
for \(1\le e\le x\), with the endpoint \(e=x+1\) handled separately by selecting all gaps.  Running the full
\(H=0,\ldots,499\) grid with \(e\le8\) gives the same displayed prefix value and bucket contributions.  On this full
grid, the largest \(e=2\) contribution relative to \(e=1\) is already far below the dominant term:
\[
\begin{array}{c|c|c}
T & \text{worst }H & \log_2(e=2)-\log_2(e=1)\\
\hline
5949 & 499 & -54.051865\\
9949 & 499 & -47.632529\\
13949 & 499 & -46.640982
\end{array}
\]
and \(e=3,\ldots,8\) are lower still.  This is still a finite numerical checkpoint, not the desired theorem: the remaining
episode proof debt is to dominate the all-\(e\) tail uniformly, or to choose a finite \(E\) and attach a certified
tail bound for \(e>E\).

The first such tail split is already adequate for the \(h\le500\) prefix.  The crude bound
\[
\sum_{e\ge e_{\min}}\binom{H+1}{e}p_{\rm term}^e
\]
controls the event that at least \(e_{\min}\) gaps terminate, using only the fact that at most \(H+1\) gaps exist after
the first active block.  This bound is much too weak for \(e_{\min}=2\): it ignores the survival/output cost left after
only a few skipped gaps, and the far \(h=500\) outer mass overwhelms it.  It is also the wrong shape for a post-prefix
tail unless \(e_{\min}\) is large enough, because it discards the live-span Chernoff cost.  However, after the exact
\(e\le8\) ledger is included, the same crude bound with \(e_{\min}=9\) gives
\[
\log_2\sum_{\substack{32\le h\le500\\1\le r\le64\\\text{three buckets}}}
A_h^{\rm RM}\,
\Pr[\text{first active occupancy }r\text{ in bucket}]\,
\sum_{e\ge9}\binom{H+1}{e}p_{\rm term}^e
=-284.004805
\]
after exact-outer reweighting.  The peak of this residual is again the far bucket at \(h=500,r=1\), but it is now far
below the explicit \(e\le8\) ledger.  Combining the explicit \(e\le8\) ledger with this \(e\ge9\) residual leaves the
current finite-prefix bound
\[
\log_2 \mu_{\le500}^{\rm fullsplit,\ exact\ outer}\le -37.383345.
\]
The tight part of this inequality is now exactly the \(h=32\) near-boundary row; the middle/far buckets, exact support
weights \(h>32\), and the crude \(e\ge9\) tail are already much lower.

For the early first-active buckets \(T>17948\), the same \(e\le8\) endpoint audit would be too expensive to promote
directly over every interior \(T\).  Instead the checked finite certificate uses a bucket-uniform survival bound.  For a
bucket \([T_-,T_+]\), remaining input weight \(H\), occupied remaining blocks \(x\), and selected skipped gaps \(e\),
write \(S(L)\) for the Chernoff survival bound with \(L\) live blocks.  Since
\[
\binom{bT}{H}\ge \binom{bT_-}{H}
\]
and, after setting \(L=T-s\),
\[
\binom{T-L+e-1}{e-1}\binom{L-e}{x-e}
\le
\binom{T_+-L+e-1}{e-1}\binom{L-e}{x-e}
\qquad (T_-\le T\le T_+),
\]
the whole bucket is bounded by replacing the skipped-gap suffix by
\[
\binom{x+1}{e}
\sum_{L=x}^{T_+}
\binom{T_+-L+e-1}{e-1}\binom{L-e}{x-e}S(L),
\]
with the endpoint \(e=x+1\) handled as \(\binom{T_+}{x}S(x)\).  This keeps the live-span Chernoff cost but removes the
need for an all-interior \(T\) audit.  The generator is
\[
\texttt{python scripts/build\_fullsplit\_early\_uniform\_survival.py --e-max 16}.
\]
Using the three early buckets \(12001\text{--}17000\), \(17001\text{--}22000\), \(22001\text{--}26819\), the resulting
row ledger gives
\[
\log_2 \mu_{32\le h\le500,\ {\rm early}}^{e\le16}\le -85.403338,
\]
with bucket contributions
\[
12001\text{--}17000:\ -85.815880,\qquad
17001\text{--}22000:\ -87.831441,\qquad
22001\text{--}26819:\ -89.394315.
\]
The peak row is \(h=32,r=1\) in the first early bucket, with inner exponent \(-108.576588\) and full term
\(-87.363392\).  The crude turnoff-count-only residual is then safe once started at \(e_{\min}=17\):
\[
\log_2 \mu_{32\le h\le500,\ {\rm early}}^{e\ge17}\le -305.738816.
\]

\begin{corollary}[Small-weight finite-prefix certificate]
\label{cor:fullsplit-small-prefix-all-first-active}
At the finite checkpoint \(N=2^{21}\), \(\delta=.09\), and the RM direct-sum outer from
Lemma~\ref{lem:rm-prefix-support-split}, the contribution of all first-active positions for supported weights
\(32\le h\le500\) satisfies
\[
\log_2\mu_{32\le h\le500}^{\rm all\ first\ active,\ exact\ outer}\le -37.278528.
\]
More explicitly, the three placement regimes obey
\[
\begin{array}{c|c}
\text{first-active regime} & \log_2\text{ contribution}\\
\hline
T<5949 & -41.113442\\
5949\le T\le17948 & -37.383345\\
T>17948 & -86.910456
\end{array}
\]
after exact RM outer reweighting.
\end{corollary}

\begin{proof}
The ultra-late line \(T<5949\) is the placement-only exact-outer prefix in
Lemma~\ref{lem:rm-prefix-support-split}.  The window line \(5949\le T\le17948\) is the log-sum of the explicit
\(e\le8\) ledger, justified over the three window buckets by
Lemma~\ref{lem:fullsplit-tiny-prefix-interiorT}, and the crude \(e\ge9\) residual
\[
\log_2\mu_{32\le h\le500,\ {\rm window}}^{e\ge9,\ {\rm exact\ outer}}\le -284.004805.
\]
The early line \(T>17948\) is the log-sum of the bucket-uniform \(e\le16\) ledger and the crude \(e\ge17\) residual:
\[
\log_2\mu_{32\le h\le500,\ {\rm early}}^{e\le16,\ {\rm exact\ outer}}\le -86.910456,
\qquad
\log_2\mu_{32\le h\le500,\ {\rm early}}^{e\ge17,\ {\rm exact\ outer}}\le -319.977808.
\]
Finally,
\[
\log_2\!\left(2^{-41.113442}+2^{-37.383345}+2^{-86.910456}\right)\le -37.278528.
\]
The verifier \(\texttt{scripts/verify\_fullsplit\_finite\_ledger.py}\) checks each displayed numerical input and the
final log-sum.
\end{proof}

For the remaining early branch \(h>500\), high-density rows expose an important bookkeeping trap.  If one reuses the
endpoint-placement interface above, the scan can look catastrophically positive: near \(h=1100001\) in the far early
bucket, endpoint placement at \(T_{\max}\) paired with the inner survival bound at \(T_{\min}=27949\) gives a false
high-density obstruction.  For \(H\) comparable with \(bT\), however, the placement mass in the bucket is concentrated
near large \(T\), exactly where the survival exponent is stronger.

The working replacement is therefore the \(T\)-paired sum
\[
\sum_{T=T_-}^{T_+}
\binom{b}{r}\binom{bT}{H}\binom{N}{h}^{-1}
\exp_2\{S(T)\},
\]
where \(S(T)\) is the fixed-\(\lambda\) survival exponent for \(T\) live blocks, with an episode multiplier attached
afterwards.  The diagnostic
\[
\texttt{python scripts/probe\_fullsplit\_early\_paired\_survival.py}
\]
keeps this pairing.  On the representative formerly bad row \(h=1100001\), far early bucket, the paired no-turnoff
sum is already
\[
\log_2 \mu_{\rm paired}^{e=0}\le -133084.867453,
\]
peaking at \(T=32767\).  Even after adding a deliberately pessimistic \(20000\)-bit episode allowance, the same row is
\[
\log_2 \mu_{\rm paired}^{\rm slack}\le -113084.867453.
\]
Sparse high-density probes with this \(20000\)-bit allowance give
\[
\begin{array}{c|c|c}
h\text{ range sampled} & \text{sample step} & \log_2\text{ paired contribution}\\
\hline
300001\text{--}1200000 & 100000 & -112940.527337\\
1200001\text{--}2097152 & 100000 & -111645.563005 .
\end{array}
\]
These are probes, not yet a certificate over every \(h\).  They show that the apparent high-density early obstruction
comes from separating endpoint placement from endpoint survival; the next theorem target is a paired
placement-survival interval lemma, with a fixed-pole episode multiplier attached to the same \(T\)-sum.
There is a second, separate bookkeeping lesson at the near-full-density edge.  If \(H\) nearly equals \(bT\), the
moderate-\(\rho\) Cauchy coefficient bound in \eqref{eq:fullsplit-e1-cauchy-envelope} can also look much too large:
it upper-bounds the exact coefficient of \(((1+z)^b-1)^x\) from the wrong side of the polynomial.  The fix is not an
extra heuristic slack term but a large-\(\rho\) (equivalently, deficit-side) pole.  At the forced endpoint
\(T=32767,H=bT=2097088,\lambda=2.3\), the same eall-ratio wrapper changes from about
\(-893425.892127\) with \(\rho=10\) to \(-1181735.473865\) with \(\rho=10^6\), matching the no-turnoff survival
exponent.  Thus the near-full edge is also a coefficient-envelope issue, not a new episode obstruction.
On the outer side, the concrete RM512 local spectrum is complement-symmetric, and hence the direct-sum outer satisfies
\(A_h^{\rm out}=A_{N-h}^{\rm out}\).  The high-\(h\) early branch should therefore be certified by the same
low/complement-weight outer table used below \(N/2\), together with the paired-\(T\) inner sum.  The current scripts
expose this as the \texttt{--outer-complement-symmetry} option; at \(h=N,r=64,T=32767\) the paired survival smoke gives
\(-1181735.473865\), while the endpoint piecewise form gives \(-1181723.239347\), differing only by the
\(\log_2(4819)\) bucket-size overcount.

The paired probe now also supports the same-\(T\) all-episode wrapper directly.  A sparse survival scan over
\(h=501,100501,\ldots,2000501,N\), with complement-symmetric outer weights and all \(1\le r\le64\), peaks at
\[
h=1000501,\qquad r=31,\qquad T=32767,
\qquad
\log_2\mu^{e=0}_{\rm paired}\le -136323.731114.
\]
Attaching the all-episode \(e\)all-ratio bound at the same \(T\)-values and scanning local \(r\)-windows near the
survival peaks gives
\[
\begin{array}{c|c|c|c}
h & r\text{-window} & \log_2\mu_{\rm paired}^{\rm eall} & \text{peak }r\\
\hline
1000501 & 25\text{--}37 & -133179.527574 & 32\\
1150501 & 30\text{--}42 & -132946.839848 & 32\\
1200501 & 32\text{--}43 & -132647.058774 & 32\\
1250501 & 33\text{--}45 & -132835.935679 & 33 .
\end{array}
\]
This is still a sampled diagnostic, not a closed interval theorem.  Its value is that the remaining early
high-density certificate no longer looks numerically tight: the proof debt is to package the paired \(T\)-sum and the
local \(r\)-window reduction into interval lemmas.
The first-block count direction also appears benign.  On the same four sampled \(h\)-values, if one freezes the
dominating endpoint \(T=32767\) and sums over all \(1\le r\le64\), the paired all-episode totals are
\[
\begin{array}{c|c}
h & \log_2\mu_{\rm paired}^{\rm eall}(T=32767,\;1\le r\le64)\\
\hline
1000501 & -133179.355669\\
1150501 & -132946.388557\\
1200501 & -132646.193762\\
1250501 & -132834.786925 .
\end{array}
\]
For the worst of these rows, \(h=1200501\), the single peak term is \(-132649.525145\) at \(r=32\), so the entire
\(r\)-sum costs only \(3.331383\) bits.  Meanwhile the single-\(r=32\), all-\(T\) sum is
\(-132649.521954\), only \(0.003191\) bits above the endpoint peak.  The helper
\texttt{sum\_fullsplit\_paired\_early\_eall.py} now packages this in the theorem-facing form: it splits the
all-episode wrapper into the \(e=0\) branch and the \(e\ge1\) fixed-pole Cauchy branch, pairs each branch with the
same \(T\)-placement term, and then sums \(T\) directly.  On the full far bucket \(22001\le g\le26819\), with all
\(1\le r\le64\), it gives
\[
\begin{array}{c|c}
h & \log_2\mu_{\rm paired}^{\rm eall}(22001\le g\le26819,\;1\le r\le64)\\
\hline
1000501 & -133179.352477\\
1150501 & -132946.385366\\
1200501 & -132646.190570\\
1250501 & -132834.783733 .
\end{array}
\]
At \(h=1200501\) this is only \(0.003192\) bits above the endpoint/full-\(r\) value, and the displayed total is
entirely dominated by the \(e\ge1\) branch; the \(e=0\) branch contributes \(-136323.731114\).  For the pole used in
this row, \(\lambda=2.3,\rho=1\), one has
\[
\log_2 G_{\lambda,\rho}=8.821856,
\]
so Lemma~\ref{lem:fullsplit-paired-endpoint-dominance} gives a uniform \(T\)-bucket overhead of at most
\[
-\log_2\left(1-\frac{1}{G_{\lambda,\rho}-1}\right)=0.003199
\]
bits for the \(e\ge1\) branch.  This turns the observed endpoint dominance into a proof-shaped bound.  The remaining
early high-density task is therefore an exact hypergeometric first-block \(r\)-sum or \(r\)-tail bound, plus an interval
bound in \(h\) for the outer/complement spectrum.

A broader endpoint scan shows why the \(h\)-interval step should not be skipped.  Scanning
\[
h=1100501,1105501,\ldots,1300501
\]
moves the sampled peak to the neighborhood of \(h=1115\)k.  A step-one zoom over \(1115101\le h\le1115201\) finds
\[
h=1115146
\]
as the sampled endpoint peak.  At this row, the full far bucket \(22001\le g\le26819\), all \(1\le r\le64\), gives
\[
\log_2\mu_{\rm paired}^{\rm eall}\le -132529.969694,
\]
with endpoint value \(-132529.972886\), \(T\)-overhead \(0.003192\) bits, and peak \(r=32\) endpoint term
\(-132533.304269\).  Thus the corrected sampled high-density peak is about \(116\) bits larger than the earlier
four-point table suggested, but still has the same structural explanation: endpoint dominance in \(T\), a small
constant first-block \(r\)-sum, and a large remaining first-moment margin.

The interval version of this argument is implemented in
\[
\texttt{scripts/certify\_fullsplit\_high\_density\_interval.py}.
\]
For a fixed outer Chernoff pole \(z\), the complement-side term
\[
\log A_{N-h}^{\rm out}(z)-\log\binom Nh-h\log_2\rho
\]
is linear in \(N-h\) plus \(-\log\binom{N}{N-h}\), hence convex on any \(h\)-interval.  The interval helper therefore
checks only the endpoints for the outer/placement part, uses Lemma~\ref{lem:fullsplit-paired-endpoint-dominance} for
the \(T\)-sum, and charges the whole first-block sum by
\[
\log_2\sum_{r=1}^{64}\binom{64}{r}\rho^r .
\]
With the RM512 complement-symmetric outer, this gives the following pessimistic complement-side certificate:
\[
\begin{array}{c|c|c}
h\text{ range} & \rho & \log_2\mu\\
\hline
1048577\text{--}1148736 & 1 & -128915.315746\\
1148737\text{--}1300000 & 1 & -124889.651755\\
1300001\text{--}1500000 & 1 & -113710.357653\\
1500001\text{--}1700000 & 1 & -35150.097366\\
1700001\text{--}1900000 & 10 & -524628.366199\\
1900001\text{--}2097089 & 10 & -894140.350713\\
2097090\text{--}2097152 & 10 & -893792.285010 .
\end{array}
\]
The combined complement-side total is \(-35150.097366\) bits.  This is much weaker than the sampled local peak, but it
is summable and closes the high-density complement branch with only two coefficient poles, \(\rho=1\) and \(\rho=10\).
The finite-ledger verifier records this endpoint reduction as
\(\texttt{complement\_high\_shape\_status=PASS}\); the smallest discrete convexity-growth certificate over the listed
intervals is \(0.276449723567\) bits.

For the post-prefix direction, a product-generating-function RM outer scan was run with \(h=501,\ldots,2000\), the
same three gap buckets and \(1\le r\le64\), and full integer \(H=0,\ldots,1999\) \(e\le1\) episode grids.  This gives
\[
\log_2 \mu_{501\le h\le2000}^{(e\le1),{\rm gf}}
\le -202.271474,
\]
with the aggregate peak already at \(h=501\), gap \(1\text{--}4000\).  Thus the exact episode branch is not developing
a new intermediate-\(h\) ridge beyond the prefix.  By contrast, the crude turnoff-count-only \(e\ge9\) residual alone is
not usable post-prefix: over the same \(501\le h\le2000\) window it becomes positive and peaks at the artificial
endpoint, solely because it ignores the survival cost after the skipped gaps.  A high-end exact profile confirms the
right target.  At \(H=1999,T=13949\),
\[
\begin{array}{c|cccc}
e & 1 & 2 & 4 & 8\\
\hline
\log_2\text{ contribution} & -2529.389774 & -2572.382804 & -2662.541447 & -2850.958022
\end{array}
\]
so the exact selected-gap terms still decay rapidly in \(e\).  The next theorem-shaped object should therefore be an
\emph{exact episode-tail ratio} for the closed-form \(e\)-summands, not a tail bound that counts turnoffs without
charging the remaining live span.  The product-GF post-prefix scanner is
\[
\texttt{python scripts/sum\_fullsplit\_episode\_outer\_gf.py}.
\]
The all-episode wrapper below now supplies this replacement on the same \(501\le h\le2000\) range.  With modes
\[
\texttt{cap,eallratio,eallratio},
\]
it gives
\[
\log_2 \mu_{501\le h\le2000}^{\rm all\mbox{-}episode}
\le -182.259739,
\]
again peaking at \(h=501,r=1\) in the near bucket.  This is weaker than the \(e\le1\) diagnostic, but it removes the
post-prefix multi-episode caveat.
The fast episode diagnostic, interpolation checker, and crude tail summation helper are
\[
\begin{gathered}
\texttt{python scripts/fast\_fullsplit\_episode\_e01.py},\\
\texttt{python scripts/check\_fullsplit\_e\_tail\_ratio.py},\\
\texttt{python scripts/check\_fullsplit\_interpolation\_slack.py},\\
\texttt{python scripts/sum\_fullsplit\_turnoff\_tail.py}.
\end{gathered}
\]

\begin{lemma}[Product bound for selected full-split terminations]
\label{lem:fullsplit-selected-gap-product}
Fix the full-codeword random-split inner.  Suppose \(p_{\rm term}\) is a uniform conditional upper bound on the
one-live-branch termination atom under the exposure convention used by the episode ledger.  Namely, after exposing any
earlier blocks, earlier scramblers and splits, the current state weight, the gap/occupancy skeleton being summed over,
and the remaining input-placement counts, but before exposing the fresh random labels that can cause the next charged
termination, the conditional probability of that charged event is at most \(p_{\rm term}\).  The charged event may be the
split atom \(q'=0\), or an exact input/state cancellation atom bundled into the same opportunity.

Then, for any fixed exposed gap/occupancy skeleton and any fixed set \({\cal S}\) of \(e\) selected gaps in the episode
decomposition,
\[
\Pr[\text{all selected opportunities in }{\cal S}\text{ terminate}\mid\text{skeleton}]
\le
p_{\rm term}^{e}.
\]
\end{lemma}

\begin{proof}
Order the selected opportunities from left to right and expose the process sequentially.  Let \({\cal F}_{j-1}\) be the
history before the \(j\)-th selected opportunity, including the outcomes of all earlier selected and unselected blocks.
On the event that the \(j\)-th opportunity is actually reached with a live state, the definition of \(p_{\rm term}\)
gives
\[
\Pr[E_j\mid {\cal F}_{j-1}]\le p_{\rm term}.
\]
If the opportunity is not reached, the joint event that all selected opportunities terminate has conditional probability
zero from that history.  Multiplying these conditional bounds by the chain rule gives the claimed \(p_{\rm term}^e\)
bound.

For the full-split construction the hypothesis is exactly the local exposure principle used above.  Conditional on a
nonzero incoming state, the fresh scrambler makes the BCH codeword uniform over nonzero codewords, and the fresh split
controls the \(q'=0\) atom independently of the past.  The exact-cancellation part is whatever local random-placement
atom has been left unexposed in the chosen ledger; the finite-prefix and global values of \(p_{\rm term}\) below are the
corresponding certified union bounds for the split atom and that cancellation atom.  The global atom is pointwise in a
fixed next-block support, while the sharper finite-prefix atom is used only with the finite ledger exposure for which it
is verified.
\end{proof}

\begin{lemma}[All-episode generating bound for the full-split inner]
\label{lem:fullsplit-all-episode-gf-bound}
Fix the full-codeword random-split inner with block size \(b\).  Let
\[
p_{\rm term}
\]
be any valid upper bound on the one-live-branch termination atom, including both the \(q'=0\) split event and any
exact input/state cancellation event being charged in the same step.  Let \(p_{j,q}\) be the nonzero BCH split law, and
for \(\lambda>0\) set
\[
M_+(\lambda):=\sum_{j,q>0} p_{j,q}e^{-\lambda j}.
\]
Suppose the first active block has already been exposed, and \(H\) input ones remain to be placed uniformly in \(T\)
later \(b\)-blocks.  Let
\[
C_{H,x}^{(b)}:=[z^H]\big((1+z)^b-1\big)^x
\]
be the number of placements of the remaining \(H\) ones using exactly \(x\) nonempty later blocks.  Then for every
\(\lambda>0\), the full selected-gap episode contribution, summed over all numbers of terminating/skipped gaps, is at
most
\begin{equation}
\label{eq:fullsplit-all-e-gf-exactcoeff}
\sum_x
\frac{C_{H,x}^{(b)}}{\binom{bT}{H}}\,
e^{\lambda d}M_+(\lambda)^T\,
[t^{T-x}]
\left(
\frac{1}{1-t}
+
\frac{p_{\rm term}}{1-t/M_+(\lambda)}
\right)^{x+1}.
\end{equation}
Equivalently, for any \(a\) with \(0<a<M_+(\lambda)\), the same contribution is bounded by
\begin{equation}
\label{eq:fullsplit-all-e-gf-cauchy}
\sum_x
\frac{C_{H,x}^{(b)}}{\binom{bT}{H}}\,
e^{\lambda d}M_+(\lambda)^T\,
a^{-(T-x)}
\left(
\frac{1}{1-a}
+
\frac{p_{\rm term}}{1-a/M_+(\lambda)}
\right)^{x+1}.
\end{equation}
\end{lemma}

\begin{proof}
Condition on exactly \(x\) nonempty later blocks.  The number of such placements is
\(\binom{T}{x}C_{H,x}^{(b)}\), so the occupancy probability is
\[
\frac{\binom{T}{x}C_{H,x}^{(b)}}{\binom{bT}{H}}.
\]
Given \(x\), the \(G=T-x\) empty blocks are distributed uniformly among the \(x+1\) gaps after active blocks, including
the terminal suffix.  Select \(e\) of these \(x+1\) gaps to be skipped by terminations.  If the selected gaps contain
total length \(s\), then the number of gap compositions with this selected total is
\[
\binom{s+e-1}{e-1}\binom{T-s-e}{x-e}
\]
with the endpoint conventions \(e=0\) and \(e=x+1\).  The factor
\(\binom{T}{x}\) in the occupancy probability cancels the denominator
\(\binom{G+x}{x}=\binom{T}{x}\) from the uniform gap-composition law.

By Lemma~\ref{lem:fullsplit-selected-gap-product}, any fixed selected set contributes at most \(p_{\rm term}^e\).  If
the selected gaps have total length \(s\), then \(T-s\) live BCH split branches remain to be charged.  For an exact
cancellation charged to a preceding gap, deleting the intervening live branches is only pessimistic in the following
Chernoff bound because \(M_+(\lambda)\le1\).  For fixed \(\lambda\),
\[
\Pr[\text{live output weight}\le d]
\le e^{\lambda d}M_+(\lambda)^{T-s}
=e^{\lambda d}M_+(\lambda)^T\left(M_+(\lambda)^{-1}\right)^s.
\]
Thus one selected gap has ordinary generating function
\((1-t/M_+(\lambda))^{-1}\) and contributes an additional factor \(p_{\rm term}\), while an unselected gap has
generating function \((1-t)^{-1}\).  Summing over all selected subsets of all sizes gives the coefficient in
\eqref{eq:fullsplit-all-e-gf-exactcoeff}.  The Cauchy bound
\([t^G]F(t)\le a^{-G}F(a)\), valid for \(0<a<M_+(\lambda)\), gives
\eqref{eq:fullsplit-all-e-gf-cauchy}.
\end{proof}

\begin{corollary}[Coefficient-free all-episode envelope]
\label{cor:fullsplit-all-episode-coeff-free}
In the setting of Lemma~\ref{lem:fullsplit-all-episode-gf-bound}, fix \(\lambda>0\), \(0<a<M_+(\lambda)\), and
\(\rho>0\), and write
\[
\Gamma_{\lambda,a}:=
\frac{1}{1-a}
+
\frac{p_{\rm term}}{1-a/M_+(\lambda)}.
\]
Let
\[
x_0(H):=
\begin{cases}
0,&H=0,\\
\lceil H/b\rceil,&H>0,
\end{cases}
\qquad
x_1(H,T):=\min\{H,T\}.
\]
Then the all-episode contribution is at most
\begin{equation}
\label{eq:fullsplit-all-e-coeff-free}
\frac{e^{\lambda d}M_+(\lambda)^T}{\binom{bT}{H}}\,
\Gamma_{\lambda,a}\,a^{-T}\rho^{-H}
\sum_{x=x_0(H)}^{x_1(H,T)}
\left(a\Gamma_{\lambda,a}\big((1+\rho)^b-1\big)\right)^x .
\end{equation}
\end{corollary}

\begin{proof}
Starting from \eqref{eq:fullsplit-all-e-gf-cauchy}, factor
\(\Gamma_{\lambda,a}a^{-T}\) from the summand.  The remaining occupancy sum is
\[
\sum_x C_{H,x}^{(b)}(a\Gamma_{\lambda,a})^x.
\]
For every \(\rho>0\),
\[
C_{H,x}^{(b)}
= [z^H]\big((1+z)^b-1\big)^x
\le
\rho^{-H}\big((1+\rho)^b-1\big)^x,
\]
and \(C_{H,x}^{(b)}=0\) outside the displayed range of \(x\).  This gives
\eqref{eq:fullsplit-all-e-coeff-free}.
\end{proof}

This corollary is the proof object that should replace the finite \(e\)-grid wherever a smooth tail envelope is needed.
It is local-code agnostic except through the split law \(p_{j,q}\) and the termination upper bound \(p_{\rm term}\), and
it interfaces with any outer code only through the usual first-moment sum over \(A_h^{\rm out}\).  In particular, RM,
BCH-like block outers, and a random-banded outer with an analytic spectrum envelope can all use the same inner episode
theorem.  The outer code supplies either exact coefficients or an envelope
\[
A_h^{\rm out}\le \mathcal A(h),
\]
and the dense+dense first moment then uses the same schematic sum
\[
\sum_h \mathcal A(h)\sum_r
\Pr[\text{first active inner block has weight }r\mid h]\,
\mathcal E(h-r,T_r,d),
\]
where \(\mathcal E\) is any certified inner envelope, such as
\eqref{eq:fullsplit-all-e-coeff-free}.  Thus the RM block outer is a concrete calibration instance, not part of the
inner theorem.

The first coefficient-free diagnostic is mixed in exactly the useful way.  With the EBCH \([128,64]\) split law,
the global theorem-facing termination atom \(p_{\rm term}=2^{-62.4078758}\), \(N=2^{21}\), \(b=64\), and \(\delta=.09\),
optimizing a coarse
\((\lambda,a,\rho)\)-grid gives
\[
\begin{array}{c|cc}
T & H=499 & H=1999\\
\hline
5949 & 0 & 0\\
9949 & -169.605380 & 0\\
13949 & -355.018537 & -278.016997 .
\end{array}
\]
The zeros mean that the probability cap \(1\) is the best value found by this smooth all-episode envelope.  Therefore
the full proof should be explicitly piecewise: short remaining spans are handled by outer-spectrum and placement
suppression, while middle/far spans use the all-episode inner envelope.  The remaining analytic task is to optimize
\((\lambda,a,\rho)\) in \eqref{eq:fullsplit-all-e-coeff-free} and prove a post-prefix summable outer/inner product
bound, rather than to extend the finite \(e\)-grid.

The corresponding summation wrapper is
\[
\texttt{python scripts/sum\_fullsplit\_piecewise\_certificate.py}.
\]
It takes a local outer spectrum, first-active placement buckets, and a per-bucket inner mode.  The current theorem-shaped
configuration is
\[
\texttt{cap,cauchy,cauchy},
\]
meaning that the near bucket uses the trivial inner probability bound \(1\), while the middle and far buckets use
Corollary~\ref{cor:fullsplit-all-episode-coeff-free}.  This is intentionally a certificate interface rather than an
RM-specific evaluator: replacing \(RM(4,9)\) by a BCH-like block outer, or by an analytic random-banded spectrum
envelope, changes only the outer generating function supplied to the first-moment sum.

The first run of this wrapper clarified the remaining proof split.  The pure smooth configuration
\[
\texttt{cap,cauchy,cauchy}
\]
is \emph{not} adequate on the small-\(H\) ridge: over \(32\le h\le80\) it gives
\[
\log_2\mu=37.549030,
\]
peaking at \(h=76,r=1\) in the far bucket, because the coefficient-free all-episode Cauchy bound is essentially capped
there.  Replacing the middle and far buckets by the audited fixed-\(e\) grids fixes this ridge:
\[
\begin{array}{c|c|c}
\text{range} & \text{inner modes} & \log_2\mu\\
\hline
32\le h\le80 & \texttt{csv,csv,csv} & -34.767371\\
32\le h\le500 & \texttt{csv,csv,csv} & -34.767174\\
501\le h\le2000 & \texttt{csv,csv,csv}\ (e\le1\text{ grids}) & -202.271474\\
501\le h\le2000 & \texttt{cap,eallratio,eallratio} & -182.259739 .
\end{array}
\]
The first two rows use the \(e\le8\) grids through \(H=499\); the sharper post-prefix diagnostic uses the \(e\le1\)
grids through \(H=1999\), while the last row is the theorem-facing all-episode replacement.  In all cases the peak is
at the left edge of the relevant range, so the finite-grid prefix is not hiding an intermediate-weight ridge.

For the tail beyond the audited grid, the right smooth object is not the all-\(e\) Cauchy envelope but a fixed-\(e=1\)
envelope.  For \(H>0\), the \(e=0\) branch is bounded by
\[
E_0(\lambda)=e^{\lambda d}M_+(\lambda)^T,
\]
while the one-skipped-gap branch admits, for every \(\rho>0\),
\begin{equation}
\label{eq:fullsplit-e1-cauchy-envelope}
E_1(\lambda,\rho)
\le
\frac{p_{\rm term}e^{\lambda d}}{\binom{bT}{H}}\,
\rho^{-H}
\sum_{x=\lceil H/b\rceil}^{\min\{H,T\}}
(x+1)
\left(
\big((1+\rho)^b-1\big)\frac{M_+(\lambda)}{1-M_+(\lambda)}
\right)^x .
\end{equation}
This comes from the exact \(e=1\) suffix identity
\[
\sum_{\ell=x}^{T}\binom{\ell-1}{x-1}M_+(\lambda)^\ell
\le
\left(\frac{M_+(\lambda)}{1-M_+(\lambda)}\right)^x
\]
and the coefficient bound
\[
C_{H,x}^{(b)}\le \rho^{-H}\big((1+\rho)^b-1\big)^x.
\]
Numerically, at the bad middle-tail point \(T=9949,H=3000\), this gives
\[
\log_2 E_{\le1}\le -2163.662004,
\]
compared with the exact fixed-\(e\) value \(-2298.278345\).  Thus it preserves the important part of the exact
episode-grid strength while avoiding a high-\(H\) coefficient table.

The same fixed-pole calculation gives the desired multi-termination tail comparison.  Fix \(\lambda\) and write
\[
M=M_+(\lambda),\qquad p=p_{\rm term}.
\]
For fixed \(x\), let \(\widehat E_{e,x}^{(\lambda)}\) denote the contribution with exactly \(x\) occupied later blocks
and \(e\) charged terminations, using the survival bound \(e^{\lambda d}M^L\) at live length \(L\).  Then
\begin{equation}
\label{eq:fullsplit-fixed-x-e-tail-ratio}
\sum_{e\ge2}\widehat E_{e,x}^{(\lambda)}
\le
\widehat E_{1,x}^{(\lambda)}
\sum_{k=1}^{x}
\frac{\binom{x+1}{k+1}}{x+1}
\binom{T-x+k}{k}
\bigl(p(1-M)\bigr)^k .
\end{equation}
Consequently, after summing over \(x\),
\begin{equation}
\label{eq:fullsplit-e-tail-ht-multiplier}
\widehat E_{\ge2}^{(\lambda)}
\le
\widehat E_{1}^{(\lambda)}\,
R_{H,T}(\lambda),
\qquad
R_{H,T}(\lambda):=
\sum_{k=1}^{\min\{H,T\}}
\frac{\binom{H}{k}\binom{T}{k}}{k+1}
\bigl(p(1-M_+(\lambda))\bigr)^k .
\end{equation}
Indeed, for \(1\le e\le x\), writing \(k=e-1\) and \(L=T-s\), the fixed-\(e\) suffix sum is bounded by
\[
\sum_{L=x}^{T}
\binom{T-L+e-1}{e-1}
\binom{L-e}{x-e}M^L
\le
\binom{T-x+e-1}{e-1}
\frac{M^x}{(1-M)^{x-e+1}}.
\]
The \(e=1\) envelope has the same \(C_{H,x}^{(b)}\), occupancy denominator, and factor
\[
p\,e^{\lambda d}(x+1)\frac{M^x}{(1-M)^x}.
\]
Dividing gives the summand in \eqref{eq:fullsplit-fixed-x-e-tail-ratio}.  The endpoint \(e=x+1\), where all gaps are
selected, gives the same \(k=x\) factor directly.  Finally,
\[
\frac{\binom{x+1}{k+1}}{x+1}=\frac{\binom{x}{k}}{k+1}\le \frac{\binom{H}{k}}{k+1},
\qquad
\binom{T-x+k}{k}\le \binom{T}{k},
\]
which gives \eqref{eq:fullsplit-e-tail-ht-multiplier}.

The preceding derivation is the theorem-facing object used by the finite ledger:

\begin{lemma}[Fixed-pole all-episode ratio envelope]
\label{lem:fullsplit-fixed-pole-eallratio}
In the setting of Lemma~\ref{lem:fullsplit-all-episode-gf-bound}, assume \(H>0\), fix \(\lambda>0\) with
\(M:=M_+(\lambda)<1\), and fix \(\rho>0\).  Put
\[
E_0(\lambda):=e^{\lambda d}M^T,
\]
and
\[
B_1(\lambda,\rho):=
\frac{p_{\rm term}e^{\lambda d}}{\binom{bT}{H}}\,
\rho^{-H}
\sum_{x=\lceil H/b\rceil}^{\min\{H,T\}}
(x+1)
\left(
\big((1+\rho)^b-1\big)\frac{M}{1-M}
\right)^x .
\]
With \(R_{H,T}(\lambda)\) defined in \eqref{eq:fullsplit-e-tail-ht-multiplier}, the full selected-gap episode
contribution is at most
\[
E_0(\lambda)+B_1(\lambda,\rho)\bigl(1+R_{H,T}(\lambda)\bigr).
\]
Consequently the inner probability contribution may be capped by the minimum of this quantity and \(1\).
\end{lemma}

\begin{proof}
The \(e=0\) branch has no charged termination and is bounded directly by \(E_0(\lambda)\).  For \(e=1\), conditioning on
the number \(x\) of occupied later blocks gives the exact suffix sum
\[
(x+1)\sum_{\ell=x}^{T}\binom{\ell-1}{x-1}M^\ell
\le
(x+1)\left(\frac{M}{1-M}\right)^x,
\]
while Cauchy's coefficient bound gives
\[
C_{H,x}^{(b)}
\le
\rho^{-H}\big((1+\rho)^b-1\big)^x .
\]
Summing over \(x\) gives \(B_1(\lambda,\rho)\).  Finally,
\eqref{eq:fullsplit-e-tail-ht-multiplier} compares the \(e\ge2\) contribution to the same fixed-\(\lambda\) \(e=1\)
contribution by the multiplier \(R_{H,T}(\lambda)\).  Adding the three branches gives the claim.
\end{proof}

\begin{lemma}[Right-endpoint dominance for paired early sums]
\label{lem:fullsplit-paired-endpoint-dominance}
Fix \(b,H,\lambda,\rho\), and write
\[
G_{\lambda,\rho}:=
\big((1+\rho)^b-1\big)\frac{M_+(\lambda)}{1-M_+(\lambda)}.
\]
Assume \(G_{\lambda,\rho}>2\).  Let \(T_0\le T_1 < H\), and put
\[
T_*:=\max\{T_0,\lceil H/b\rceil\},\qquad
S_T(H;\lambda,\rho):=
\sum_{x=\lceil H/b\rceil}^{T}(x+1)G_{\lambda,\rho}^x .
\]
Then
\begin{equation}
\label{eq:fullsplit-paired-endpoint-dominance}
\sum_{T=T_*}^{T_1}
S_T(H;\lambda,\rho)\bigl(1+R_{H,T}(\lambda)\bigr)
\le
\frac{
S_{T_1}(H;\lambda,\rho)\bigl(1+R_{H,T_1}(\lambda)\bigr)
}{
1-(G_{\lambda,\rho}-1)^{-1}
}.
\end{equation}
\end{lemma}

\begin{proof}
For \(T\ge T_*\),
\[
S_{T-1}
\le
\sum_{x=\lceil H/b\rceil}^{T-1}T\,G_{\lambda,\rho}^x
\le
\frac{T\,G_{\lambda,\rho}^{T}}{G_{\lambda,\rho}-1},
\qquad
S_T\ge (T+1)G_{\lambda,\rho}^{T}.
\]
Thus \(S_{T-1}/S_T\le (G_{\lambda,\rho}-1)^{-1}\).  Iterating this adjacent-ratio bound gives a geometric tail
from the right endpoint.  The multiplier \(R_{H,T}(\lambda)\) in
\eqref{eq:fullsplit-e-tail-ht-multiplier} is nondecreasing in \(T\), since every summand contains
\(\binom{T}{k}\).  Replacing all \(1+R_{H,T}\) by the endpoint value \(1+R_{H,T_1}\) and summing the geometric series
gives \eqref{eq:fullsplit-paired-endpoint-dominance}.
\end{proof}

The multiplier is tiny in the finite regimes that currently matter.  The audit helper
\[
\texttt{python scripts/check\_fullsplit\_e\_tail\_multiplier.py}
\]
reports, for example, the following conservative \(R_{H,T}\) exponents:
\[
\begin{array}{c|c|c|c}
T & H & \lambda & \log_2 R_{H,T}(\lambda)\\
\hline
5949 & 10000 & .05 & -39.399491\\
9949 & 10000 & .05 & -38.657489\\
13949 & 10000 & .05 & -38.169806 .
\end{array}
\]
The sharper fixed-\(x\) multiplier is another \(1.5\)--\(2.7\) bits lower in these samples.  Thus the \(e\ge2\) tail
should add a negligible multiplicative factor to the fixed-\(\lambda\) \(e=1\) envelope, provided the first-moment
summation keeps the same \(\lambda\) when applying \eqref{eq:fullsplit-e-tail-ht-multiplier}.
The verifier now also audits the numerical evaluation of \(R_{H,T}\).  The implementation truncates the positive
\(k\)-sum only after adding a geometric upper bound for the omitted suffix, using the adjacent-ratio formula
\[
\frac{a_{k+1}}{a_k}
=
\frac{(H-k)(T-k)}{(k+1)(k+2)}\,p(1-M_+(\lambda)).
\]
For the current feasible \(e\)all-ratio range \(T\le32767\), this ratio is already below
\(2^{-28.992971}\) after the first term, even under the crude bound \(1-M_+(\lambda)\le1\).  The finite-ledger
manifest records exact finite-sum samples against this truncated evaluator as
\(\texttt{eallratio\_tail\_status=PASS}\).
The same manifest now also includes a small implementation audit of the whole fixed-pole branch formula.  For six finite
\((T,H,\lambda,\rho)\) samples, the verifier computes the selected-gap Chernoff sum directly from exact integer
coefficients of \(((1+z)^b-1)^x\) and exact skipped-gap suffix sums, then checks that
\eqref{eq:fullsplit-e1-cauchy-envelope} with the same-\(\lambda\) multiplier
\eqref{eq:fullsplit-e-tail-ht-multiplier} dominates it.  It also compares the production
\(\texttt{eallratio}\) evaluator against an independent fixed-pole formula evaluator, with maximum discrepancy
\(2.9\cdot10^{-14}\) bits in the current manifest.  This is an implementation audit of the algebra; the proof object is
still the displayed coefficient, suffix, and multiplier inequalities above.

The wrapper now has a theorem-facing mode
\[
\texttt{cap,eallratio,eallratio},
\]
where the near bucket is still charged only by placement, while the middle and far buckets use
Lemma~\ref{lem:fullsplit-fixed-pole-eallratio}, namely \eqref{eq:fullsplit-e1-cauchy-envelope} together with the
same-\(\lambda\) multiplier \eqref{eq:fullsplit-e-tail-ht-multiplier}.  Thus this mode is an all-episode upper bound,
not merely an \(e\le1\) diagnostic.
The finite wrapper grid now also includes large \(\rho\)-poles up to \(10^4\).  These poles are mostly irrelevant in
the ordinary sparse/intermediate occupancy range, but are essential near \(H=bT\), where they implement the same
deficit-side control one would get by rewriting the occupancy coefficient in terms of missing coordinates.

For fast high-weight exploration we also use the rigorous endpoint placement bound
\[
\sum_{g=g_0}^{g_1}\binom{b(T_0+g-1)}{H}
\le
(g_1-g_0+1)\binom{b(T_0+g_1-1)}{H}.
\]
This loses at most \(\log_2 4000<12\) bits per bucket in the current three-bucket split, but avoids millions of
independent binomial evaluations.  The finite-ledger verifier also samples exact bucket sums at the feasibility
boundaries and records the endpoint dominance check; the largest advertised bucket-size loss is \(11.965784\) bits.
The matching inner endpoint must be handled with some care: if the left endpoint \(T_{\min}\) has \(H>bT_{\min}\),
then early gaps in the bucket are impossible but later gaps are not.  The corrected working rule is therefore to use
the earliest feasible block count
\[
T_{\rm eff}(H;T_{\min},T_{\max})
:=\max\{T_{\min},\lceil H/b\rceil\},
\]
discarding the bucket only when \(T_{\rm eff}>T_{\max}\).  Turning this into a theorem requires the mild monotonicity
lemma that the fixed-pole envelope is nonincreasing in \(T\) for \(T\ge \lceil H/b\rceil\); otherwise one should split
the bucket further.

Here is the exact monotonicity condition now used for audits.  For fixed \((\lambda,\rho)\), put
\[
M=M_+(\lambda),\qquad
G_{\lambda,\rho}:=
\big((1+\rho)^b-1\big)\frac{M}{1-M},
\]
and define
\[
S_T(H;\lambda,\rho):=
\sum_{x=\lceil H/b\rceil}^{\min\{H,T\}}(x+1)G_{\lambda,\rho}^x.
\]
The \(e\ge1\) part of the fixed-pole envelope is monotone from \(T\) to \(T+1\) whenever
\[
\log_2\frac{\binom{b(T+1)}{H}}{\binom{bT}{H}}
\ge
\log_2\frac{S_{T+1}(H;\lambda,\rho)}{S_T(H;\lambda,\rho)}
+
\log_2\frac{1+R_{H,T+1}(\lambda)}{1+R_{H,T}(\lambda)}.
\]
The \(e=0\) part is already monotone because it is proportional to \(M^T\) and \(M<1\).  The audit script
\[
\texttt{scripts/check\_fullsplit\_T\_monotonicity.py}
\]
checks this one-step inequality on the interval certificate.

There is a simple sufficient endpoint reduction.  Write
\[
D_{H,T}:=\frac{\binom{b(T+1)}{H}}{\binom{bT}{H}}
=\prod_{i=1}^{b}\frac{bT+i}{bT+i-H}.
\]
Each factor is increasing in \(H\) and decreasing in \(T\).  Thus the no-new-occupancy case \(T\ge H\), where
\(S_{T+1}=S_T\), is worst at the smallest \(H\) in the interval and the largest \(T\) in the bucket.  In the
new-occupancy case \(T<H\), the smallest feasible \(H\) is \(\max\{H_{\min},T+1\}\), and
\[
\frac{S_{T+1}}{S_T}
=
1+\frac{(T+2)G_{\lambda,\rho}^{T+1}}{S_T}
\le
1+\frac{T+2}{T+1}G_{\lambda,\rho},
\]
since \(S_T\ge (T+1)G_{\lambda,\rho}^T\).  Finally \(R_{H,T}(\lambda)\) is increasing in both \(H\) and \(T\), so
\[
\log_2\frac{1+R_{H,T+1}(\lambda)}{1+R_{H,T}(\lambda)}
\le
\log_2\bigl(1+R_{H_{\max},T_{\max}}(\lambda)\bigr).
\]
The command
\[
\texttt{python scripts/check\_fullsplit\_T\_monotonicity.py --sufficient-reduction}
\]
checks these sufficient endpoint inequalities.  With the interval poles below, every middle/far bucket has positive
slack:
\[
\begin{array}{c|c|c|c}
h\text{ range} & (\lambda,\rho) & \text{middle slack} & \text{far slack}\\
\hline
2001\text{--}7858 & (.02,.01) & .200576 & .155836\\
7859\text{--}20550 & (.05,.03) & .154987 & .154991\\
20551\text{--}75000 & (.2,.1) & 1.129733 & .647617\\
75001\text{--}250000 & (.5,.3) & 4.010933 & 2.144089\\
250001\text{--}350000 & (1.2,.5) & 30.056467 & 22.381864\\
350001\text{--}650000 & (1.2,1) & 21.636936 & 9.238901\\
650001\text{--}725000 & (1.2,2) & 58.501371 & 15.287874\\
725001\text{--}950000 & (1.2,3) & 66.069027 & 3.773620\\
950001\text{--}1148736 & (1.2,3) & \infty & 73.689638 .
\end{array}
\]
This is now a finite endpoint certificate rather than a dense scan over \(H,T\): the denominator growth beats the
new occupancy term and the \(R_{H,T}\) correction by a visible margin in every interval.
The finite-ledger verifier runs this sufficient-reduction audit by default; its manifest records the worst row as
\(7859\text{--}20550\) in the middle bucket, with slack \(0.154987\) bits.

With the endpoint placement bound, the all-episode wrapper gives the following dense checked ranges:
\[
\begin{array}{c|c|c}
\text{range} & \log_2\mu & \text{dominant bucket}\\
\hline
2001\le h\le5000 & -750.899709 & \text{near, cap}\\
5001\le h\le10000 & -1916.794192 & \text{near, cap}\\
10001\le h\le20000 & -3622.467998 & \text{middle, }e\text{-all}\\
20001\le h\le50000 & -6021.761195 & \text{middle, }e\text{-all}.
\end{array}
\]
Using the corrected \(T_{\rm eff}\) rule and the current adaptive pole grid, a coarse sanity scan over
\(250000\le h\le1000000\) in steps of \(25000\) gives \(\log_2\mu=-78531.366918\), peaking at \(h=250000\).
This is only a sanity scan: the dense interval table below is the relevant certificate, since a coarse step can miss an
interior peak.  The corrected scan replaces the earlier misleading high-tail ``impossible bucket'' observation, which
came from pairing endpoint placement with the bucket's left endpoint \(T\).

Even before optimizing the high tail, a fixed-pole interval certificate already works through the feasible far-bucket
tail \(h=1148736\):
\[
\begin{array}{c|c|c}
h\text{ range} & (\lambda,\rho) & \log_2\mu\\
\hline
2001\text{--}7858 & (.02,.01) & -586.597103\\
7859\text{--}20550 & (.05,.03) & -2655.944828\\
20551\text{--}75000 & (.2,.1) & -6025.574047\\
75001\text{--}250000 & (.5,.3) & -12844.868837\\
250001\text{--}350000 & (1.2,.5) & -28476.640541\\
350001\text{--}400000 & (1.2,1) & -81555.594538\\
400001\text{--}450000 & (1.2,1) & -91421.772986\\
450001\text{--}550000 & (1.2,1) & -82081.961328\\
550001\text{--}650000 & (1.2,1) & -1518.493461\\
650001\text{--}725000 & (1.2,2) & -116328.133946\\
725001\text{--}750000 & (1.2,3) & -169887.902921\\
750001\text{--}850000 & (1.2,3) & -134523.583646\\
850001\text{--}950000 & (1.2,3) & -117534.843530\\
950001\text{--}1050000 & (1.2,3) & -285052.550668\\
1050001\text{--}1148736 & (1.2,3) & -562860.716638 .
\end{array}
\]
The displayed high-interval rows are the raw finite-prefix interval constants with a checked \(1.484774\)-bit
adjustment for replacing the finite-prefix termination atom by the global Bernstein termination atom.  The helper
\(\texttt{verify\_fullsplit\_high\_interval\_adjustment.py}\) computes the required adjustment as
\(1.484773404861\) bits, including a negligible tail-multiplier allowance.  The ledger helper
\(\texttt{verify\_fullsplit\_finite\_ledger.py}\) can print the exact fixed-pole regeneration command for each interval
and can spot-recompute selected intervals using \(\texttt{--check-high-intervals}\).
The tight row remains the upper transition interval \(550001\text{--}650000\), not the multi-episode correction itself.
Above this transition the feasible far-bucket tail quickly falls away; the later peaks are either at the left endpoint or
near \(h=862232\), and all are far below the target first-moment level.
The cutoff is the feasibility edge of the current late-window split: the far bucket has
\(T_{\max}=5949+12000-1=17948\), so \(H\le 64T_{\max}=1148672\), and hence no row with \(r\le64\) exists once
\(h>1148736\).
The complement-side interval table above is the replacement for this old feasible-cutoff limitation on the high-density
early branch.  It uses the paired \(T\)-sum directly at the true far endpoint \(T=32767\), the complement symmetry
\(A_h^{\rm out}=A_{N-h}^{\rm out}\), and fixed-\(z\) endpoint bounds over \(h\)-intervals.  The finite ledger records
this as
\[
\log_2\mu_{1048577\le h\le N}^{\rm complement\ high}\le -35150.097366.
\]

It remains to include the post-prefix part of the ultra-late branch \(T<5949\).  Here no inner episode estimate is
needed: all \(h\) nonzero coordinates must lie in the final \(m=64\cdot5949=380736\) inner coordinates.  For a fixed
outer Cauchy pole \(z\),
\[
A_h^{\rm out}\frac{\binom{m}{h}}{\binom Nh}
\le
\left(W_{\rm loc}(z)^{4096}-1\right)z^{-h}\frac{\binom{m}{h}}{\binom Nh}.
\]
The logarithm of the right-hand side has adjacent ratio \((m-h)/(N-h)z^{-1}\), hence is concave in \(h\).  Therefore
each interval is certified by checking its endpoints and the critical point
\[
h_*=\frac{m-zN}{1-z}.
\]
With \(z=0.39605985943459426\), the checked intervals
\[
\begin{array}{c}
501\text{--}2000,\quad 2001\text{--}7858,\quad 7859\text{--}20550,\\
20551\text{--}75000,\quad 75001\text{--}150000,\quad 150001\text{--}250000,\\
250001\text{--}380736
\end{array}
\]
have union-bound total
\[
\log_2\mu_{501\le h\le380736,\ T<5949}^{\rm placement}\le -545.690526.
\]
The finite-ledger verifier also checks the adjacent-ratio shape reduction.  For this pole all seven intervals are
maximized at their left endpoint.  The verifier flag \(\texttt{--print-late-postprefix-commands}\) prints the
regeneration commands.

\begin{corollary}[Checked post-prefix ledger rows]
\label{cor:fullsplit-postprefix-checked-ledger}
The checked \(h\ge501\) row families currently included in the finite ledger have union-bound total
\[
\log_2\mu_{\rm checked,\ h\ge501}^{\rm fullsplit}\le -182.259739.
\]
These rows are
\[
\begin{array}{c|c}
\text{row family} & \log_2\text{ contribution}\\
\hline
501\le h\le2000\ \text{all-episode wrapper} & -182.259739\\
501\le h\le380736,\ T<5949\ \text{ultra-late placement rows} & -545.690526\\
501\le h\le75000,\ T>17948\ \text{early endpoint rows} & -380.829185\\
75001\le h\le350000,\ T>17948\ \text{accelerated endpoint rows} & -1056.187188\\
350001\le h\le N/2,\ T>17948\ \text{paired accelerated rows} & -98386.449256\\
2001\le h\le1148736\ \text{fixed-pole intervals} & -586.597103\\
1048577\le h\le N\ \text{complement-high intervals} & -35150.097366 .
\end{array}
\]
\end{corollary}

\begin{proof}
The first row is the tracked \(501\le h\le2000\) all-episode h-summary artifact for the late window.  The second row is
the ultra-late placement cap for \(h>500\).  Since \(64\cdot5949=380736\), no ultra-late placement row exists above
\(h=380736\).  The third row is the checked early endpoint table for \(T>17948\) through \(h=75000\).  The fourth row is
the vectorized endpoint extension for \(75001\le h\le350000\), and the fifth row is the paired-\(T\) extension for
\(350001\le h\le N/2\).  The verifier checks that these accelerated early intervals form a contiguous cover
\(75001,\ldots,N/2\), split into seven endpoint-helper rows on \(75001,\ldots,350000\) and three paired-helper rows on
\(350001,\ldots,N/2\), with no gap or overlap.  The verifier also checks that the ultra-late placement rows cover
\(501,\ldots,380736=64\cdot5949\).  The sixth row is the fixed-pole interval table through the feasible far-bucket
cutoff \(h=1148736\); its verifier entry checks a contiguous \(2001,\ldots,1148736\) cover and the common
\(1.484774\)-bit turnoff adjustment.  It also checks that the feasible-min rule
\(T_{\rm eff}=\max(T_{\min},\lceil H/64\rceil)\) gives far-bucket cutoff \(64\cdot17948+64=1148736\), and records the
endpoint gap-sum audit and sufficient \(T\)-monotonicity audit, with monotonicity minimum slack \(0.154987\) bits.  The
seventh row is the complement-side high-density table, checked as a contiguous \(1048577,\ldots,N\) cover with the
convex endpoint reduction above.  The last two row families overlap on \(1048577\le h\le1148736\); the ledger records
this overlap and deliberately keeps it as harmless union-bound overcount.  Therefore the combined contribution is
bounded by
\[
\begin{aligned}
&\log_2\!\left(2^{-182.259739}+2^{-545.690526}+2^{-380.829185}+2^{-1056.187188}\right.\\
&\qquad\left.+2^{-98386.449256}+2^{-586.597103}+2^{-35150.097366}\right)\le -182.259739.
\end{aligned}
\]
The verifier prints and checks this as
\(\texttt{h501\_plus\_checked\_rows\_log2=-182.259739}\).
\end{proof}

\begin{corollary}[Current finite first-moment checkpoint]
\label{cor:fullsplit-current-finite-checkpoint}
Let \(\mu_{\rm checked}(.09)\) denote the first-moment contribution of the row families covered by
Corollaries~\ref{cor:fullsplit-small-prefix-all-first-active} and~\ref{cor:fullsplit-postprefix-checked-ledger}.
Then
\[
\log_2 \mu_{\rm checked}(.09)\le -37.278528.
\]
\end{corollary}

\begin{proof}
Corollary~\ref{cor:fullsplit-small-prefix-all-first-active} gives
\[
\log_2\mu_{32\le h\le500}^{\rm all\ first\ active,\ exact\ outer}\le -37.278528,
\]
and Corollary~\ref{cor:fullsplit-postprefix-checked-ledger} gives
\[
\log_2\mu_{\rm checked,\ h\ge501}^{\rm fullsplit}\le -182.259739.
\]
Thus
\[
\log_2\!\left(2^{-37.278528}+2^{-182.259739}\right)\le -37.278528.
\]
The verifier prints and checks the same log-sum under the
\(\texttt{current\_checked\_ledger}\) row, with value \(-37.278528\).
\end{proof}

\begin{theorem}[Checked finite RM/full-split first-moment certificate]
\label{thm:fullsplit-rm-finite-certificate-009}
Consider the finite checkpoint \(N=2^{21}\), \(b=64\), and \(\delta=.09\), so that
\[
d=\lfloor .09N\rfloor=188743 .
\]
Use the RM direct-sum outer from Lemma~\ref{lem:rm-prefix-support-split} and the full-split dense inner certified
above.  Let \(Z_d\) be the number of nonzero concatenated codewords whose final weight is at most \(d\), with the
expectation taken over the random choices of the inner construction.  The checked finite ledger gives
\[
\mathbb E Z_d\le 2^{-37.278528}.
\]
Consequently
\[
\Pr[d_{\min}\le d]\le 2^{-37.278528}.
\]
In particular, with positive probability a sampled code has \(d_{\min}\ge188744\), and hence relative distance
\[
\frac{d_{\min}}{N}\ge \frac{188744}{2^{21}}>.09 .
\]
\end{theorem}

\begin{proof}
Lemma~\ref{lem:rm-prefix-support-split} verifies that the RM direct-sum outer has no positive weights below \(32\).
Corollary~\ref{cor:fullsplit-small-prefix-all-first-active} covers all supported weights \(32\le h\le500\) and all
first-active positions by splitting them into the ultra-late, window, and early regimes.  Corollary~
\ref{cor:fullsplit-postprefix-checked-ledger} covers the remaining outer weights \(h\ge501\): the ultra-late rows cover
all feasible \(T<5949\) cases, the all-episode and fixed-pole interval rows cover the late window, the early endpoint
and paired rows cover \(T>17948\) through \(N/2\), and the complement-high rows cover \(h>N/2\).  The overlap between
the fixed-pole and complement-high rows is retained only as union-bound overcount.

Thus every nonzero outer word is included in the row families defining \(\mu_{\rm checked}(.09)\).  Corollary~
\ref{cor:fullsplit-current-finite-checkpoint} gives \(\log_2\mu_{\rm checked}(.09)\le -37.278528\), which is exactly
the displayed bound on \(\mathbb E Z_d\).  The probability statement is Markov's inequality applied to the nonnegative
integer random variable \(Z_d\).
\end{proof}

\subsection{Spectrum-model outer projections}
\label{subsec:fullsplit-outer-spectrum-projections}

The proved checkpoint above uses the exact RM \([512,256,32]\) local spectrum.  As a finite-\(N\) design diagnostic, we
also compared two possible stronger outer floors while leaving the full-split EBCH inner ledger unchanged.  This is not
a theorem upgrade: the BCH rows below use a spectrum model, not a certified local spectrum or a rigorous low-weight
envelope.  The comparison script is
\[
\texttt{python scripts/compare\_outer\_modes\_fullsplit.py --delta 0.09}.
\]
It writes the CSV and Markdown reports
\[
\begin{gathered}
\texttt{scripts/outer\_mode\_comparison\_delta009.csv},\\
\texttt{scripts/outer\_mode\_comparison\_delta009.md}.
\end{gathered}
\]
The RM row reproduces the checked ledger value
\(-37.278528\), while the BCH rows are marked heuristic in both outputs.  The modeled local spectrum is
\[
A_w=0\quad(0<w<d_0),\qquad
A_w\approx 2^{k_0-n_0}\binom{n_0}{w}\quad(w\ge d_0),
\]
with optional sensitivity inflation of \(+10,+20,+40\) bits on the low-weight window
\(d_0\le w\le d_0+32\).

At \(N=2^{21}\), \(\delta=.09\), and \(d=188743\), the baseline projection rows are:
\[
\begin{array}{c|c|c|c|c|c|c}
\text{outer mode} & \text{status} & [n_0,k_0,d_0] & \text{blocks} &
\log_2\mu & \text{margin bits} & \text{dominant }h\\
\hline
\text{RM exact} & \text{proved} & [512,256,32] & 4096 & -37.278528 & 37.278528 & 32\\
\text{BCH-like} & \text{heuristic} & [256,128,38] & 8192 & -37.089134 & 37.089134 & 59\\
\text{extended BCH-like} & \text{heuristic} & [512,256,\ge62] & 4096 & -68.171304 & 68.171304 & 119 .
\end{array}
\]
All three baseline rows are still dominated by the first-active near-window row \(r=1\), \(1\le\mathrm{gap}\le4000\).
The rough work proxy \(B n_0\log_2 n_0\) is \(18874368\) for the \(512\)-bit rows and \(16777216\) for the \(256\)-bit
row, so the shorter-block projection has similar or slightly lower nominal outer work.

The sensitivity check is the warning sign.  The \([256,128,38]\) model remains near the RM margin with no inflation and
has \(27.167072\) bits of margin under \(+10\) bits, but the projected \(501\le h\le2000\) row becomes dominant and
positive under \(+20\) bits.  The \([512,256,\ge62]\) model is much more robust in this diagnostic: its margins under
\(+10,+20,+40\) bits are \(65.802609\), \(56.111333\), and \(36.111659\) bits, respectively.  Thus the projection
suggests that a real \(512\)-bit BCH-like outer spectrum could be a meaningful constant-factor upgrade, whereas the
\(256\)-bit option depends sensitively on the actual low-weight spectrum.

This subsection should be read only as a spectrum-model projection.  To make any BCH row theorem-facing, we still need
either the exact local weight distribution or a proved low-weight envelope, plus a recomputation of the \(h>2000\) tail
under that same outer model.  The report records the external pointers used for the projection notes: the CodeTables
\([256,128]\) construction record, the Kusaka/Okayama weight-distribution index, and the BCH designed-distance bound.

The finite ledger is now checked by
\[
\texttt{python scripts/verify\_fullsplit\_finite\_ledger.py}.
\]
The same verifier writes a machine-readable snapshot by passing \(\texttt{--write-manifest-json}\) with target
\(\texttt{scripts/fullsplit\_finite\_ledger\_manifest.json}\).
That manifest records the checked row families, fixed poles, recomputation commands, and log-sum totals under schema
version \(v1\).
It first checks the tracked interior-\(T\) audit artifact for the \(32\le h\le500\), \(e\le8\) window prefix, then combines
that explicit CSV ledger with the crude \(e\ge9\) window residual, the bucket-uniform early \(h\le500\) ledger
\(e\le16\), the crude early \(e\ge17\) residual, the all-episode \(501\le h\le2000\) wrapper row, the checked
ultra-late post-prefix placement rows, the checked early endpoint rows through \(h=75000\), the accelerated early rows
through \(h=N/2\), and the fixed-pole interval tables above, including the complement-side high-density manifest.
Reweighting the \(h\le500\) pieces by exact
RM direct-sum outer coefficients gives the reported combined finite late-window contribution
\[
\log_2\mu_{\rm finite}^{\rm fullsplit,\ exact\ outer}\le -37.383345,
\]
dominated by the tiny-prefix row \(h=32,r=1\) in the near bucket; the exact-outer \(e\ge9\) prefix residual is already
\(-284.004805\) bits, the exact-outer early \(h\le500\) contribution is \(-86.910456\) bits, and the entire \(h\ge501\)
checked ledger is below \(-182.259739\) bits.  Including the checked exact-outer ultra-late placement prefix, the
all-position \(32\le h\le500\) value is
\[
\log_2\mu_{32\le h\le500}^{\rm all\ first\ active\ positions,\ exact\ outer}\le -37.278528.
\]

The per-\(h\) summary output also gives a compact active-pole map through \(h=50000\).  The dominant row is:
\[
\begin{array}{c|c|c|c}
h\text{ range} & \text{dominant bucket} & (\lambda,\rho) & \text{dominant }r\\
\hline
2001\text{--}7357 & \text{near, cap} & - & 1\\
7358\text{--}7858 & \text{middle} & (.02,.01) & 1\\
7859\text{--}10748 & \text{middle} & (.05,.03) & 1\\
10749\text{--}13781 & \text{near, cap} & - & 1\\
13782\text{--}17314 & \text{far} & (.05,.03) & 1\\
17315\text{--}20550 & \text{middle} & (.05,.03) & 1\text{--}4\\
20551\text{--}22707 & \text{middle} & (.2,.1) & 4\\
22708\text{--}50000 & \text{near, cap} & - & 2\text{--}5 .
\end{array}
\]
Thus, through the main dense checked band, only four all-episode pole choices ever dominate:
\[
(.02,.01),\quad (.05,.03),\quad (.2,.1),\quad (.5,.3)
\]
with larger occupancy poles appearing only in the corrected high-tail scan.  This makes the next theorem target
plausibly finite: prove the \(T_{\rm eff}\)-monotonicity lemma, then prove a small interval certificate for these pole
choices using the endpoint placement bound and the outer generating-function envelope.

The important lesson from these scans is methodological.  The \(e\ge2\) tail is no longer the blocker once it is
charged against the same fixed pole as the \(e=1\) branch.  The remaining proof obligation is to replace the finite
adaptive pole grid by a clean analytic split in \(H/T\): small \(H\) wants small \(\lambda,\rho\), intermediate \(H\)
wants \(\lambda\approx .05\text{--}.2\), and high \(H\) wants \(\lambda\approx .5\text{--}1.2\) with larger occupancy
poles.  This is a standard Chernoff-envelope bookkeeping problem, rather than a new episode-combinatorics problem.

% ============================================================
